\chapter{Утверждения}\label{sec:statements}\label{chap:statements}
% Обновить примеры схем / R1CS: описать их языки и как выглядят наивные доказательства

% https://people.cs.georgetown.edu/jthaler/ProofsArgsAndZK.pdf
% https://docs.zkproof.org/reference.pdf
Как мы видели в неформальном введении, SNARK — это краткий неинтерактивный аргумент знания (succinct non-interactive argument of knowledge), где доказательство знания подтверждает корректность утверждений (statements) вроде ``Доказывающий знает разложение на простые множители данного числа'' или ``Доказывающий знает прообраз данного значения SHA2 digest'' и тому подобное. Однако человекочитаемые утверждения такого рода неформальны и не очень полезны с формальной точки зрения.

В этой главе мы поэтому рассмотрим более подробно способы формализации утверждений математически строгими методами, полезными для разработки SNARK. Мы начнём с введения формальных языков как способа правильного определения утверждений (\secname{} \ref{sec:formal-languages}). Для подробного введения в формальные языки см., например, \cite{moll-2012}. Затем мы рассмотрим алгебраические схемы (algebraic circuits) и \concept{системы ограничений ранга 1} [R1CS] как два особенно полезных способа определения утверждений в определённых формальных языках (\secname{} \ref{sec:statement-representations}). \concept{Системы ограничений ранга 1} и алгебраические схемы введены, например, в приложении E к \cite{sasson-2013}.

Правильное проектирование утверждений (statement design) должно иметь высокий приоритет при разработке SNARK, поскольку непреднамеренные истинные утверждения могут привести к потенциально серьёзным и почти не обнаруживаемым уязвимостям безопасности в приложениях SNARK.

\section{Формальные языки}\label{sec:formal-languages}

Формальные языки предоставляют теоретический фундамент, в рамках которого утверждения могут быть сформулированы логически строгим образом, а доказательство корректности любого данного утверждения может быть реализовано вычислением слов в этом языке.

Можно возразить, что понимание формальных языков не очень важно при разработке SNARK и проектировании связанных утверждений, но термины из этой области исследований являются стандартным жаргоном во многих статьях о доказательствах с нулевым разглашением. Поэтому мы считаем, что хотя бы краткое введение в формальные языки и то, как они вписываются в картину разработки SNARK, полезно — в основном для того, чтобы дать разработчикам лучшую интуицию о том, где всё это находится в более широкой картине логического ландшафта. Кроме того, формальные языки дают лучшее понимание того, чем на самом деле является формальное доказательство утверждения.

Грубо говоря, формальный язык (или просто язык) — это множество слов. Слова, в свою очередь, являются строками букв, взятых из некоторого алфавита, и сформированных в соответствии с определяющими правилами языка.

Будем более точны. Пусть $\Sigma$ — произвольное множество и $\Sigma^*$ — множество всех \term{строк} (strings) конечной длины $<x_1,\ldots,x_n>$ элементов $x_j$ из $\Sigma$, включая пустую строку $<>\in \Sigma^*$. Тогда \term{язык} $L$, в своём наиболее общем определении, является подмножеством множества всех конечных строк $\Sigma^*$. В этом контексте множество $\Sigma$ называется \term{алфавитом} (alphabet) языка $L$, элементы из $\Sigma$ называются \term{буквами} (letters), а элементы из $L$ называются \term{словами} (words).
Если существуют правила, определяющие, какие строки из $\Sigma^*$ принадлежат языку, а какие нет, эти правила называются \term{грамматикой} (grammar) языка. Если $L_1$ и $L_2$ — два формальных языка над одним и тем же алфавитом, мы называем $L_1$ и $L_2$ \term{эквивалентными} (equivalent), если они состоят из одного и того же множества слов.\footnote{Более подробное объяснение этого определения можно найти, например, в \secname{} 1.2 книги \cite{moll-2012}.}

\begin{comment}
While the term language might suggest a deeper relation to well-known \term{natural languages} like English, formal languages and natural languages differ in many ways. The following two examples  provide some intuition about formal languages and their differences to natural languages: 
\begin{example}[The English language] The English alphabet is given by the following set:
\begin{multline}
\Sigma_{English} = \{A, a, B, b, C, c, D, d, E, e, F, f, G, g, H, h, I, i, J, j, K, k, L, l, M, m, N, n, O,\\ o, P, p, Q, q, R, r, S, s, T, t, U, u, V, v, W, w, X, x, Y, y, Z, z\}
\end{multline}
\smecomb{The natural language English does not define words by abstract rules, but by something we might call historical human convention as a word in the language English is distinguished from any other string from the English alphabet by human convention not by formal rules. In that regard, strings like ``tea'' or ``eat'' are words in English, but ``aet'' and ``tae'' are not, because people agree on that. So while the set of all English words is a formal language $L_{English}\subset \Sigma_{English}$ there is no formal grammar that defines what a word in English is and what not.}{SB: none of this is correct, needs to be rewritten or cut}\sme{alternative example: blick vs. bnick as possible words}

It should also be recognized that in such a simple definition, the ''grammar'' of the natural language English is different from the grammar of the formal language English. The grammar of natural English defines how sentences are generated from words, while sentences have no meaning in our simple formal language English. 

In order to express natural English sentences in a formal English language, one has to extend the alphabet by a space character together with the set of all punctuation symbols,  and then include the grammar of the natural language English into the formal language English. It should be noted, however, that in such a formal language, sentences like "I drink tea, but I don't eat meat." from the natural language English are actually words in the formal language English.\sme{SB: since we're not using sentences, I'd cut this}
\end{example}
\end{comment}
\begin{example}[Чередующиеся двоичные строки]
\label{ex:alternating_strings_lang}
Чтобы рассмотреть очень простой формальный язык с почти тривиальной грамматикой, возьмём множество из двух букв, $0$ и $1$, в качестве нашего алфавита $\Sigma$:
$$
\Sigma = \{0,1\}
$$
Кроме того, мы используем грамматику, утверждающую, что правильное слово должно состоять из чередующихся двоичных букв произвольной длины, включая пустую строку. Ассоциированный язык $L_{alt}$ является множеством всех конечных двоичных строк, где за $1$ должно следовать $0$ и наоборот. Так, например, $<1,0,1,0,1,0,1,0,1>\in L_{alt}$ является словом в этом языке, как и $<0>\in L_{alt}$ или пустое слово $<>\in L_{alt}$. Однако двоичная строка $<1,0,1,0,1,0,1,1,1>\in \{0,1\}^*$ не является правильным словом, так как она нарушает грамматику $L_{alt}$, поскольку последние 3 буквы равны $1$. Кроме того, строка $<0,A,0,A,0,A,0>$ не является правильным словом, так как не все её буквы из алфавита $\Sigma$.
\end{example}

\subsection{Функции решения}
\label{sec:decision_function}
 Наше предыдущее определение формальных языков очень общее и не охватывает многие подклассы языков, известные в литературе. Однако в контексте разработки SNARK языки обычно определяются как \term{задачи разрешения} (decision problems), где так называемое \term{разрешающее отношение} (deciding relation) $R\subset \Sigma^*$ решает, является ли данная строка $x\in \Sigma^*$ словом в языке или нет. Если $x\in R$, то $x$ является словом в ассоциированном языке $L_R$, а если $x\notin R$, то нет. \href{https://en.wikipedia.org/wiki/Relation_(mathematics)}{Отношение} $R$ поэтому обобщает грамматику языка $L_R$.

К сожалению, в некоторой литературе по системам доказательств $x\in R$ часто записывается как $R(x)$, что вводит в заблуждение, поскольку, вообще говоря, $R$ не является функцией, а отношением в $\Sigma^*$. Для ясности мы поэтому примем другую точку зрения и будем работать с тем, что мы можем назвать \term{функцией решения} (decision function):
\begin{equation}\label{eq:decision-function}
R: \Sigma^* \to \{true, false\}
\end{equation}
Функции решения определяют, является ли строка $x\in \Sigma^*$ элементом языка или нет. В случае, если задана функция решения, сам ассоциированный язык можно записать как множество всех строк, которые решаются функцией $R$:
\begin{equation}
L_R := \{x\in \Sigma^*\;|\; R(x)=true\}
\end{equation}
В контексте формальных языков и задач разрешения \term{утверждение} (statement) $S$ — это утверждение о том, что язык $L$ содержит слово $x$, то есть утверждение утверждает, что существует некоторое $x\in L$. Конструктивное \term{доказательство} (proof) утверждения $S$ даётся некоторой строкой $P\in \Sigma^*$, и такое доказательство \term{верифицируется} (verified) проверкой, равно ли $R(P)=true$. В этом случае $P$ называется \term{экземпляром} (instance) утверждения $S$.

\begin{example}[Alternating Binary strings] Чтобы рассмотреть очень простой формальный язык с функцией решения, рассмотрим язык $L_{alt}$ из \examplename{} \ref{ex:alternating_strings_lang}. Попытавшись записать грамматику этого языка более формально, мы можем определить следующую функцию решения:
$$
R: \{0,1\}^* \to \{true,false\}\;;\; <x_0,x_1,\ldots,x_n> \mapsto 
\begin{cases}
true & x_{j-1} \neq x_{j} \text{ для всех } 1\leq j \leq n \\
false & \text{ иначе}
\end{cases}
$$
Мы можем использовать эту функцию, чтобы решить, является ли данная двоичная строка словом в $L_{alt}$ или нет. Некоторые примеры приведены ниже:
$$
\begin{array}{l}
R(<1,0,1>)=true\\ 
R(<0>)=true\\
R(<>)=true\\ 
R(<1,1>)=false
\end{array} 
$$
Учитывая язык $L_{alt}$, имеет смысл утверждать следующее утверждение: ``Существует чередующаяся строка.'' Один из способов доказать это утверждение \href{https://en.wikipedia.org/wiki/Constructive_proof}{конструктивно} — предоставить фактический экземпляр, то есть привести пример чередующейся строки, например $x = <1,0,1>$. Конструирование строки $<1,0,1>$ поэтому доказывает утверждение ``Существует чередующаяся строка.'', поскольку можно проверить, что $R(<1,0,1>)=true$.
\end{example}
\begin{example}[Язык программирования]Языки программирования — очень важный класс формальных языков. Для этих языков алфавит обычно является (подмножеством) таблицы ASCII, а грамматика определяется правилами компилятора языка программирования. Слова, в свою очередь, — это правильно написанные компьютерные программы, которые компилятор принимает. Компилятор поэтому можно интерпретировать как функцию решения.

Чтобы привести необычный пример, достаточно странный, чтобы подчеркнуть суть, рассмотрим язык программирования \href{https://en.wikipedia.org/wiki/Malbolge}{Malbolge}. Этот язык был специально разработан, чтобы быть почти невозможным в использовании, и написание программ на этом языке — сложная задача. Интересным утверждением является утверждение: ``Существует компьютерная программа на Malbolge''. Как оказалось, доказать это утверждение конструктивно, то есть предоставить пример экземпляра такой программы, непростая задача: потребовалось два года после появления Malbolge, чтобы написать программу, которую компилятор принимает. Так что в течение двух лет никто не мог доказать утверждение конструктивно.

Чтобы рассмотреть высокоуровневое описание Malbolge более формально, мы пишем $L_{\text{Malbolge}}$ для языка, который использует таблицу ASCII в качестве своего алфавита, а его слова — это строки ASCII-букв, которые компилятор Malbolge принимает. Доказательство утверждения ``Существует компьютерная программа на Malbolge'' эквивалентно задаче нахождения некоторого слова $x\in L_{\text{Malbolge}}$. Строка в \eqref{malbolge-string} ниже является примером такого доказательства, поскольку она принимается компилятором Malbolge, который компилирует её в исполняемый бинарник, выводящий ``Hello, World.'' \sme{add reference}. В этом примере компилятор Malbolge поэтому служит процессом верификации.

\begin{multline}\label{malbolge-string}
\scriptstyle (=<':9876Z4321UT.-Q+*)M'\&\%\$H"!~\}|Bzy?=|\{z]KwZY44Eq0/
\\ 
\scriptstyle hKs\_dG5[m\_BA\{?-Y;;Vb'rR5431M\}/.zHGwEDCBA@98\backslash 6543W10/.R,+O<
\end{multline}
\end{example}
\begin{example}[Пустой язык] Чтобы увидеть, что не каждый язык имеет хотя бы одно слово, рассмотрим алфавит $\Sigma = \Z_6$, где $\Z_6$ — кольцо модульной арифметики 6, полученное в \examplename{} \ref{def_residue_ring_z_6}. Различая множество $\Z_6^*$ всех элементов в модульной арифметике 6, имеющих мультипликативные обратные, от множества $(\Z_6)^*$ всех конечных строк над алфавитом $\Z_6$, мы определяем следующую функцию решения:\sme{SB: mention that n=1 refers to the nr. of strings?}
\begin{equation}
R_{\emptyset} : (\Z_{6})^* \to \{true, false\}\;;\;
<x_1,\ldots,x_n> \mapsto
\begin{cases}
true & n=1 \text{ и } x\cdot x = 2\\
false & \text{ иначе}
\end{cases}
\end{equation}
Мы пишем $L_\emptyset$ для ассоциированного языка. Как мы видим из таблицы умножения $\Z_6$ в \examplename{} \ref{def_residue_ring_z_6}, кольцо $\Z_6$ не содержит ни одного элемента $x$, такого что $x\cdot x =2$, что означает $R_{\emptyset}(<x_1,\ldots,x_n>)=false$ для всех строк $<x_1,\ldots,x_n>\in \Sigma^*$. Язык поэтому не содержит ни одного слова. Доказательство утверждения ``Существует слово в $L_\emptyset$'' конструктивным предоставлением экземпляра поэтому невозможно: верификация никогда не проверит ни одну строку.
\end{example}
\begin{example}[3-факторизация]\label{ex:3-factorization} Мы будем многократно использовать следующий простой пример на протяжении всей книги. Задача состоит в разработке SNARK, доказывающего знание трёх множителей элемента конечного поля $\F_{13}$. С точки зрения приложений в этом примере нет ничего особенно полезного, однако это самый простой пример, дающий нетривиальный SNARK в некоторых из наиболее распространённых систем доказательств с нулевым разглашением.

Формализуя высокоуровневое описание, мы используем $\Sigma := \F_{13}$ в качестве базового алфавита этой задачи и определяем язык $L_{3.fac}$ как состоящий из тех строк элементов поля из $\F_{13}$, которые содержат ровно $4$ буквы $x_1,x_2,x_3,x_4$, удовлетворяющие уравнению $x_1\cdot x_2\cdot x_3 =x_4$.

Так, например, строка $<2, 12, 4, 5>$ является словом в $L_{3.fac}$, в то время как ни $<2, 12, 11>$, ни $<2, 12, 4, 7>$, ни $<2, 12, 7, UPS>$ не являются словами в $L_{3.fac}$, так как они не удовлетворяют грамматике или не определены над алфавитом $\F_{13}$.

Различая множество $\F_{13}^*$ всех элементов в мультипликативной группе $\F_{13}$ от множества $(\F_{13})^*$ всех конечных строк над алфавитом $\F_{13}$, мы можем описать язык $L_{3.fac}$ более формально, введя функцию решения:
\begin{equation}
R_{3.fac} : (\F_{13})^* \to \{true, false\}\;;\;
<x_1,\ldots,x_n> \mapsto
\begin{cases}
true & n=4 \text{ и } x_1\cdot x_2 \cdot x_3 = x_4\\
false & \text{ иначе}
\end{cases}
\end{equation}
Определив язык $L_{3.fac}$, имеет смысл утверждать утверждение ``Существует слово в $L_{3.fac}$''. Как $L_{3.fac}$ спроектирован, это утверждение эквивалентно утверждению ``Существуют четыре элемента $x_1,x_2,x_3,x_4$ из конечного поля $\F_{13}$, такие что уравнение $x_1\cdot x_2\cdot x_3 =x_4$ выполняется.''

Доказательство корректности этого утверждения конструктивно означает фактически найти некоторые конкретные элементы поля, удовлетворяющие функции решения $R_{3.fac}$, например $x_1= 2$, $x_2 =12$, $x_3=4$ и $x_4 = 5$. Строка $<2,12,4,5>$ поэтому является конструктивным доказательством утверждения, что $L_{3.fac}$ содержит слова, а вычисление $R_{3.fac}(<2,12,4,5>)=true$ является верификацией этого доказательства. Напротив, строка $<2, 12, 4, 7>$ не является доказательством утверждения, поскольку проверка $R_{3.fac}(<2,12,4,7>)=false$ не верифицирует доказательство.
\end{example}
\begin{example}[Принадлежность \curvename{Tiny-jubjub}]\label{ex:tiny-jubjub} В одном из наших главных примеров мы выводим SNARK, доказывающий, что пара $(x,y)$ элементов поля из $\F_{13}$ является точкой на кривой tiny-{jubjub} в её форме скрученного Эдвардса (twisted Edwards), полученной в \examplename{} \ref{ex:TJJ13-twisted-edwards}, уравнение \ref{eq:TJJ13-twisted-edwards}.

На первом шаге мы определяем язык таким образом, что точки на кривой \curvename{Tiny-jubjub} находятся во взаимно однозначном соответствии со словами в этом языке.

Поскольку кривая \curvename{Tiny-jubjub} является эллиптической кривой над полем $\F_{13}$, мы выбираем алфавит $\Sigma = \F_{13}$. В этом случае множество $(\F_{13})^*$ состоит из всех конечных строк элементов поля из $\F_{13}$. Чтобы определить грамматику, вспомним из \eqref{TJJ13-twisted-edwards}, что точка на кривой \curvename{Tiny-jubjub} — это пара $(x,y)$ элементов поля, такая что $3\cdot x^2 + y^2 = 1+ 8\cdot x^2\cdot y^2$. Мы можем использовать это уравнение для вывода следующей функции решения:\sme{SB: mention that $x=x_1$ and $y=x_2$?}
$$
R_{tiny.jj} : (\F_{13})^* \to \{true, false\}\;;\;
<x_1,\ldots,x_n> \mapsto
\begin{cases}
true & n=2 \text{ и } 3\cdot x_1^2 + x_2^2 = 1+ 8\cdot x_1^2\cdot x_2^2\\
false & \text{ иначе}
\end{cases}
$$
Ассоциированный язык $L_{tiny.jj}$ тогда задаётся как множество всех строк из $(\F_{13})^*$, которые отображаются на $true$ функцией $R_{tiny.jj}$:
$$
L_{tiny.jj} = \{<x_1,\ldots,x_n>\in (\F_{13})^*\;|\; R_{tiny.jj}(<x_1,\ldots,x_n>)=true \}
$$
Мы можем утверждать утверждение ``Существует слово в $L_{tiny.jj}$''. Поскольку $L_{tiny.jj}$ определён $R_{tiny.jj}$, это утверждение эквивалентно утверждению ``Кривая \curvename{Tiny-jubjub} в её форме скрученного Эдвардса имеет точку кривой.'' 

Конструктивным доказательством этого утверждения является строка $<x,y>$ элементов поля из $\F_{13}$, удовлетворяющая уравнению скрученного Эдвардса. Пример \eqref{TJJ13-twisted-edwards} поэтому подразумевает, что строка $<11,6>$ является конструктивным доказательством, а вычисление $R_{tiny.jj}(<11,6>)=true$ верифицирует доказательство. Напротив, строка $<1,1>$ не является доказательством утверждения, поскольку вычисление $R_{tiny.jj}(<1,1>)=false$ не верифицирует доказательство.
\end{example}
\begin{exercise}
\label{ex:decision_function_1} Определите функцию решения таким образом, чтобы ассоциированный язык $L_{Exercise_1}$ состоял из всех решений уравнения $5x + 4 = 28 + 2x$ над $\F_{13}$. Предоставьте конструктивное доказательство утверждения: ``Существует слово в $L_{Exercise_1}$'', и верифицируйте доказательство.
\end{exercise}
\begin{exercise} Рассмотрим модульную арифметику 6 ($\Z_6$) из \examplename{} \ref{def_residue_ring_z_6}, алфавит $\Sigma = \Z_6$ и следующую функцию решения:\sme{SB: call this $R_{\Z_6}$/ $L_{\Z_6}$ instead?}
\begin{equation*}
R_{example\_\ref{def_residue_ring_z_6}} : \Sigma^*\to \{true, false\}\;;\;
<x_1,\ldots,x_n> \mapsto
\begin{cases}
true & n=1 \text{ и } 3\cdot x_1 + 3 = 0\\
false & \text{ иначе}
\end{cases}
\end{equation*}
Вычислите все слова в ассоциированном языке $L_{example\_\ref{def_residue_ring_z_6}}$, предоставьте конструктивное доказательство утверждения ``Существует слово в $L_{example\_\ref{def_residue_ring_z_6}}$'' и верифицируйте доказательство.
\end{exercise}

\subsection{Экземпляр и свидетель}
% https://www.claymath.org/sites/default/files/pvsnp.pdf
Как мы видели в предыдущем параграфе, утверждения (statements) задают утверждения о принадлежности формальным языкам, а экземпляры (instances) служат конструктивными доказательствами этих утверждений. Однако в контексте систем доказательств с \term{нулевым разглашением} (zero-knowledge) наше понятие конструктивных доказательств уточняется таким образом, что становится возможным скрыть части экземпляра доказательства и при этом иметь возможность доказать утверждение. В этом контексте поэтому необходимо разделить доказательство на не скрытую, публичную часть, называемую \term{экземпляром} (instance), и скрытую, приватную часть, называемую \term{свидетелем} (witness).

Чтобы учесть это разделение экземпляра доказательства на часть экземпляра и часть свидетеля, наше предыдущее определение формальных языков нуждается в уточнении. Вместо одного алфавита уточнённое определение рассматривает два алфавита $\Sigma_I$ и $\Sigma_W$ и функцию решения, определённую следующим образом:
\begin{equation}
R: \Sigma_I^* \times \Sigma_W^* \to \{true, false\}\;;\; (i\,;w) \mapsto R(i\,;w)
\end{equation}
Слова поэтому являются строками $(i\,;w)\in \Sigma_I^* \times \Sigma_W^*$ с $R(i\,;w)=true$. Уточнённое определение различает входы $i\in \Sigma_I$ и входы $w\in \Sigma_W$. Вход $i$ называется \term{экземпляром} (instance), а вход $w$ называется \term{свидетелем} (witness) $R$.

Если задана функция решения, ассоциированный язык определяется как множество всех строк из базовых алфавитов, которые верифицируются функцией решения:
\begin{equation}
\label{def:refined_language}
L_R := \{(i\,;w)\in \Sigma_I^* \times \Sigma_W^* \;|\; R(i\,;w)=true\}
\end{equation}
В этом уточнённом контексте \term{утверждение} (statement) $S$ — это утверждение о том, что, учитывая экземпляр $i\in\Sigma_I^*$, существует свидетель $w\in \Sigma_W^*$, такой что язык $L$ содержит слово $(i\,;w)$. Конструктивное \term{доказательство} (proof) утверждения $S$ даётся некоторой строкой $P=(i\,; w) \in \Sigma_I^* \times \Sigma_W^*$, и доказательство \term{верифицируется} (verified) проверкой $R(P)=true$.

На этом этапе важно отметить, что, хотя конструктивные доказательства в языках, различающих экземпляр и свидетель (как в \defname{} \ref{def:refined_language}), не выглядят очень по-разному от конструктивных доказательств в языках, которые мы видели в \secname{} \ref{sec:decision_function}, при наличии некоторого экземпляра существуют системы доказательств, способные доказать утверждение (по крайней мере, с высокой вероятностью), не раскрывая ничего о свидетеле, как мы увидим в \chaptname{} \ref{chapter:zk-protocols}. В этом смысле свидетель часто называют \term{приватным входом} (private input), а экземпляр — \term{публичным входом} (public input).

Стоит понимать разницу между утверждениями, определёнными в \secname{} \ref{sec:decision_function}, и уточнённым понятием утверждений из этого раздела. В то время как утверждения в смысле предыдущего раздела можно рассматривать как утверждения о принадлежности, утверждения в уточнённом определении можно рассматривать как утверждения о знании, где доказывающий утверждает знание свидетеля для данного экземпляра.
%For a more detailed discussion on this topic see [XXX\sme{update reference} sec 1.4]

\begin{example}[SHA256 --- Знание прообраза] Один из наиболее распространённых примеров в контексте систем доказательств с нулевым разглашением — это \term{доказательство знания прообраза} (knowledge-of-a-preimage proof) для некоторой криптографической \term{хеш-функции} (hash function) вроде $SHA256$, где дано публично известное значение \term{дайджеста} (digest) $SHA256$, и задача состоит в том, чтобы доказать знание прообраза для этого дайджеста под функцией $SHA256$, не раскрывая этот прообраз.

Чтобы понять эту проблему подробно, мы должны ввести язык, способный описать проблему знания прообраза таким образом, что утверждение ``Учитывая дайджест $i$, существует прообраз $w$, такой что $SHA256(w)=i$'' становится утверждением в этом языке. Поскольку $SHA256$ — это функция, отображающая двоичные строки произвольной длины на двоичные строки длины $256$:
$$
SHA256: \{0,1\}^* \to \{0,1\}^{256}
$$
Поскольку мы хотим доказать знание прообразов, мы должны рассматривать двоичные строки размера $256$ как экземпляры, а двоичные строки произвольной длины — как свидетелей.

Подходящий алфавит $\Sigma_I$ для множества всех экземпляров и подходящий алфавит $\Sigma_W$ для множества всех свидетелей поэтому задаётся множеством $\{0,1\}$. Правильная функция решения задаётся следующим образом:
\begin{multline*}
R_{SHA256} : \{0,1\}^* \times \{0,1\}^* \to \{true, false\}\;;\\
(i;w) \mapsto
\begin{cases}
true & |i|=256,\; i = SHA256(w)\\
false & \text{ иначе}
\end{cases}
\end{multline*}
Мы пишем $L_{SHA256}$ для ассоциированного языка и отмечаем, что он состоит из слов, являющихся строками $(i\,;w)$, так что экземпляр $i$ является образом $SHA256$ от свидетеля $w$.

Учитывая некоторый экземпляр $i\in \{0,1\}^{256}$, утверждение в $L_{SHA256}$ — это утверждение ``Учитывая дайджест $i$, существует прообраз $w$, такой что $SHA256(w)=i$'', что и является сутью проблемы знания прообраза. Конструктивное доказательство этого утверждения поэтому даётся прообразом $w$ для дайджеста $i$, а верификация доказательства достигается проверкой, что $SHA256(w)=i$.
\end{example}
\begin{example}[3-factorization]
\label{ex:L-3fac-zk}
 Чтобы дать ещё одну интуицию о последствиях уточнённых языков, рассмотрим снова $L_{3.fac}$ из \examplename{} \ref{ex:3-factorization}. Как мы видели, конструктивное доказательство в $L_{3.fac}$ даётся $4$ элементами поля $x_1$, $x_2$, $x_3$ и $x_4$ из $\F_{13}$, так что произведение первых трёх элементов равно $4$-му элементу в модульной арифметике 13.

Разделяя слова из $L_{3.fac}$ на части экземпляра и свидетеля, мы можем переформулировать проблему и ввести различные уровни утверждений о знании. Например, мы могли бы переформулировать утверждение о принадлежности $L_{3.fac}$ в утверждение, где все множители $x_1$, $x_2$, $x_3$ являются свидетелями, а только произведение $x_4$ — экземпляром. Утверждение для этой переформулировки тогда выражается утверждением: ``Учитывая элемент поля экземпляра $x_4$, существуют три свидетеля-множителя $x_4$''. Предполагая некоторый экземпляр $x_4$, конструктивное доказательство для ассоциированного утверждения о знании предоставляется любой строкой $(x_1,x_2,x_3)$, такой что $x_1\cdot x_2\cdot x_3= x_4$.

Мы можем формализовать этот новый язык, который мы можем назвать $L_{3.fac\_zk}$, определив следующую функцию решения:
\begin{multline*}
\label{R3-fac-zk}
R_{3.fac\_zk} : (\F_{13})^* \times (\F_{13})^* \to \{true, false\}\;;\\
(<i_1,\ldots,i_n>;<w_1,\ldots, w_m>) \mapsto
\begin{cases}
true & n=1,\; m=3,\; i_1 = w_1\cdot w_2 \cdot w_3\\
false & \text{ иначе}
\end{cases}
\end{multline*}
Ассоциированный язык $L_{3.fac\_zk}$ определяется всеми строками из $(\F_{13})^* \times (\F_{13})^*$, которые отображаются на $true$ под функцией решения $R_{3.fac\_zk}$.

Рассматривая различие, которое мы сделали между частью экземпляра и частью свидетеля в $L_{3.fac\_zk}$, можно спросить, почему мы выбрали множители $x_1$, $x_2$ и $x_3$ в качестве свидетеля, а произведение $x_4$ — в качестве экземпляра, а не другую комбинацию? Это был произвольный выбор в примере. Любая другая комбинация экземпляра и свидетеля была бы одинаково корректной. Например, можно было бы объявить все переменные свидетелем или объявить все переменные экземпляром. Фактический выбор определяется только приложением.
\end{example}
\begin{example}[Кривая \curvename{Tiny-jubjub}]
\label{ex:tiny-jubjub-instance-witness}
 Рассмотрим язык $L_{tiny.jj}$ из \examplename{} \ref{ex:tiny-jubjub}. Как мы видели, конструктивное доказательство в $L_{tiny.jj}$ даётся парой $(x_1,x_2)$ элементов поля из $\F_{13}$, так что эта пара является точкой кривой \curvename{Tiny-jubjub} в её представлении Эдвардса.

Мы рассматриваем разумное разделение слов из $L_{tiny.jj}$ на части экземпляра и свидетеля. Два очевидных выбора — либо выбрать обе координаты $x_1$ и $x_2$ как входы экземпляра, либо выбрать обе координаты $x_1$ и $x_2$ как входы свидетеля.\sme{why not 1 witness and 1 input?} 

В случае, когда обе координаты являются экземплярами, мы определяем грамматику ассоциированного языка, вводя следующую функцию решения:\sme{why are we using capital I and W here?}
\begin{multline*}
R_{tiny.jj.1} : (\F_{13})^*\times (\F_{13})^* \to \{true, false\}\;;\\
(<I_1,\ldots,I_n>;<W_1,\ldots,W_m>) \mapsto
\begin{cases}
true & n=2,\;m=0 \text{ и } 3\cdot I_1^2 + I_2^2 = 1+ 8\cdot I_1^2\cdot I_2^2\\
false & \text{ иначе}
\end{cases}
\end{multline*}
Язык $L_{tiny.jj.1}$ определяется как множество всех строк из $(\F_{13})^*\times (\F_{13})^*$, которые отображаются на $true$ функцией $R_{tiny.jj.1}$.

В случае, когда обе координаты являются входами свидетеля, мы определяем грамматику ассоциированного уточнённого языка, вводя следующую функцию решения:
\begin{multline*}
R_{tiny.jj\_zk} : (\F_{13})^*\times (\F_{13})^* \to \{true, false\}\;;\\
(<I_1,\ldots,I_n>;<W_1,\ldots,W_m>) \mapsto
\begin{cases}
true & n=0,\;m=2 \text{ и } 3\cdot W_1^2 + W_2^2 = 1+ 8\cdot W_1^2\cdot W_2^2\\
false & \text{ иначе}
\end{cases}
\end{multline*}
Язык $L_{tiny.jj\_zk}$ определяется как множество всех строк из $(\F_{13})^*\times (\F_{13})^*$, которые отображаются на $true$ функцией $R_{tiny.jj\_zk}$.
\end{example}\sme{pros and cons of $L_{tiny.jj.1}$ vs. $L_{tiny.jj\_zk}$?}


\begin{exercise} Рассмотрим модульную арифметику 6 $\Z_6$ из \examplename{} \ref{def_residue_ring_z_6} в качестве алфавитов $\Sigma_I$ и $\Sigma_W$, и следующую функцию решения:
\begin{multline*}
R_{linear} : \Sigma^* \times \Sigma^* \to \{true, false\}\;;\\
(i;w) \mapsto
\begin{cases}
true & |i|=3 \text{ и } |w|=1 \text{ и } i_1\cdot w_1 + i_2 = i_3\\
false & \text{ иначе}
\end{cases}
\end{multline*}
Какой из следующих экземпляров $(i_1,i_2,i_3)$ имеет доказательство знания в $L_{linear}$?\tbdsm{make this look nicer}
$$ 
\begin{array}{ccc}
(3,3,0) , & (2,1,0), & (4,4,2)
\end{array}
$$
\end{exercise}
\begin{exercise}[Сложение Эдвардса на \curvename{Tiny-jubjub}]
\label{ex:TJJ-addition-lang} Рассмотрим кривую \curvename{Tiny-jubjub} вместе с её законом сложения скрученного Эдвардса из \examplename{} \ref{TJJ13}. Определите алфавит экземпляра $\Sigma_I$, алфавит свидетеля $\Sigma_W$ и функцию решения $R_{add}$ с ассоциированным языком $L_{add}$, так чтобы строка $(i\,;w)\in \Sigma_I^* \times \Sigma_W^*$ была словом в $L_{add}$ тогда и только тогда, когда $i$ является парой точек кривой на кривой \curvename{Tiny-jubjub} в форме Эдвардса, а $w$ — суммой этих точек кривой.

Выберите некоторый экземпляр $i\in \Sigma_I^*$, предоставьте конструктивное доказательство утверждения ``Существует свидетель $w\in \Sigma_W^*$, такой что $(i\,;w)$ является словом в $L_{add}$'', и верифицируйте это доказательство. Затем найдите некоторый экземпляр $i\in \Sigma_I^*$, такой что $i$ не имеет доказательства знания в $L_{add}$.
\end{exercise}

\subsection{Модульность}\label{modularity} 
С точки зрения разработчика, часто полезно конструировать сложные утверждения и представляющие их языки из простых. В контексте систем доказательств с нулевым разглашением эти простые строительные блоки часто называются \term{гаджетами} (gadgets), а библиотеки гаджетов обычно содержат представления атомарных типов вроде булевых (boolean), целых чисел, различных \term{хеш-функций} (hash functions), криптографии на эллиптических кривых и многого другого (см. \chaptname{} \ref{chap:circuit-compilers}). Чтобы синтезировать утверждения, разработчики затем комбинируют предопределённые гаджеты в сложную логику. Мы называем способность комбинировать утверждения в более сложные утверждения \term{модульностью} (modularity).

Чтобы понять концепцию модульности на уровне формальных языков, определяемых функциями решения, нам нужно рассмотреть \term{пересечение} (intersection) двух языков, которое существует всякий раз, когда оба языка определены над одним и тем же алфавитом. В этом случае пересечение — это язык, состоящий из строк, которые являются словами в обоих языках.

Будем более точны. Пусть $L_1$ и $L_2$ — два языка, определённые над одними и теми же алфавитами экземпляра и свидетеля $\Sigma_I$ и $\Sigma_W$. Пересечение $L_1 \cap L_2$ языков $L_1$ и $L_2$ определяется следующим образом:
\begin{equation}
L_1 \cap L_2 := \{x\;|\; x\in L_1 \text{ и } x\in L_2\}
\end{equation} 
Если оба языка определены функциями решения $R_1$ и $R_2$, следующая функция является функцией решения для языка пересечения $L_1 \cap L_2$:
\begin{equation}
R_{L_1 \cap L_2}: \Sigma_I^* \times \Sigma_W^* \to \{true, false\}\;;\;
(i,w) \mapsto R_1(i;w) \text{ и } R_2(i;w)
\end{equation}

Таким образом, пересечение двух языков, основанных на функциях решения, также является языком, основанным на функции решения. Это важно с точки зрения реализации: оно позволяет нам конструировать сложные функции решения, их языки и ассоциированные утверждения из простых строительных блоков. Учитывая публично известный экземпляр $I\in \Sigma_I^*$, утверждение в языке пересечения утверждает знание свидетеля, удовлетворяющего всем отношениям одновременно.

\section{Представления утверждений}\label{sec:statement-representations}
Как мы видели в предыдущем разделе, формальные языки и их определения с помощью функций решения — мощный инструмент для описания утверждений формально строгим образом.

Однако с точки зрения существующих систем доказательств с нулевым разглашением не все способы фактического представления функций решения одинаково полезны. В зависимости от системы доказательства некоторые более подходят, чем другие. В этом разделе мы опишем два наиболее распространённых способа представления функций решения и их утверждений.
\subsection{Системы квадратичных ограничений ранга 1}
\label{sec:R1CS}
Хотя функции решения могут быть выражены различными способами, многие современные системы доказательства требуют, чтобы функция решения была выражена в виде системы квадратичных уравнений над конечным полем. Это верно в частности для основанных на спаривании (pairing-based) систем доказательства, подобных тем, которые мы описываем в \chaptname{}  \ref{chapter:zk-protocols}, поскольку в этих случаях возможно разделить экземпляр и свидетеля, а затем проверить решения этих уравнений ``в показателе'' (in the exponent) групп криптографии, допускающих спаривание.

В этом разделе мы более подробно рассмотрим особый тип квадратичных уравнений, называемых \term{системами квадратичных ограничений ранга 1} (Rank-1 (quadratic) Constraint Systems, R1CS), которые являются общим стандартом в системах доказательства с нулевым разглашением (см. приложение E к \cite{sasson-2013}). Мы начнём с общего введения в эти системы ограничений, а затем рассмотрим их связь с формальными языками. После этого мы рассмотрим распространённый способ вычисления решений этих систем.

\subsubsection{Представление R1CS} Чтобы понять, что такое \term{системы квадратичных ограничений ранга 1} (Rank-1 (quadratic) Constraint Systems, R1CS) подробно, пусть $\F$ — поле, $n$, $m$ и $k\in\N$ — три числа, и $a_j^i$, $b_j^i$ и $c_j^i\in\F$ — константы из $\F$ для каждого индекса $0\leq j \leq n+m$ и $1\leq i \leq k$. Тогда \concept{система ограничений ранга 1} определяется как следующее множество из $k$ уравнений:\tbdsm{font size too small} 

\begin{definition}[\deftitle{Rank-1 (quadratic) Constraint System}]\label{R1CS}
\begin{equation}\label{eq:R1CS}
\begin{split}
\scriptstyle\left(a^1_0 + \sum_{j=1}^n a^1_j \cdot I_j + \sum_{j=1}^m a^1_{n+j} \cdot W_j  \right) \cdot 
\left(b^1_0 + \sum_{j=1}^n b^1_j \cdot I_j + \sum_{j=1}^m b^1_{n+j} \cdot W_j  \right) &=\, 
\scriptstyle c^1_0 + \sum_{j=1}^n c^1_j \cdot I_j + \sum_{j=1}^m c^1_{n+j} \cdot W_j\\
       & \vdots\\
\scriptstyle\left(a^k_0 + \sum_{j=1}^n a^k_j \cdot I_j + \sum_{j=1}^m a^k_{n+j} \cdot W_j  \right) \cdot 
\left(b^k_0 + \sum_{j=1}^n b^k_j \cdot I_j + \sum_{j=1}^m b^k_{n+j} \cdot W_j  \right) &=\, 
\scriptstyle c^k_0 + \sum_{j=1}^n c^k_j \cdot I_j + \sum_{j=1}^m c^k_{n+j} \cdot W_j       
\end{split}
\end{equation}
\end{definition}

В \concept{системе ограничений ранга 1} параметр $k$ называется \term{числом ограничений} (number of constraints), и каждое уравнение называется \term{ограничением} (constraint). Если пара строк элементов поля $(<I_1,\ldots, I_n>; <W_1,\ldots,W_m>)$ удовлетворяет этим уравнениям, $<I_1,\ldots, I_n>$ называется \term{экземпляром} (instance), а $<W_1,\ldots,W_m>$ — \term{свидетелем} (witness) системы.%
\footnote{Представление \concept{систем ограничений ранга 1} может быть упрощено с использованием нотации векторов и матриц, которая абстрагирует индексы. Действительно, если $x= (1,I,W)\in \F^{1+n+m}$ — $(n+m+1)$-мерный вектор, $A$, $B$, $C$ — $(n+m+1)\times k$-мерные матрицы и $\odot$ — произведение Шура/Адамара (Schur/Hadamard product), то R1CS может быть записана следующим образом:
$$
Ax \odot Bx = Cx
$$
Однако, поскольку мы не вводили векторные пространства и матричное исчисление в книге, мы используем \ref{eq:R1CS} в качестве определяющего уравнения для \concept{систем ограничений ранга 1}. Мы лишь отметили матричную нотацию, потому что она иногда используется в литературе.}

Можно показать, что любое ограниченное вычисление может быть выражено как \concept{система ограничений ранга 1}. R1CS поэтому является универсальной моделью для ограниченных вычислений. Мы получим некоторое представление о распространённых подходах к компиляции ограниченных вычислений в \concept{системы ограничений ранга 1} в \chaptname{} \ref{chap:circuit-compilers}. 

В общем смысле идея \concept{системы ограничений ранга 1} состоит в том, чтобы отслеживать все значения, которые может принимать любая переменная во время вычисления, и связывать отношения между всеми этими переменными, которые вытекают из самого вычисления. Как только отношения между всеми шагами компьютерной программы ограничены, выполнение программы затем принудительно вычисляется точно ожидаемым образом без каких-либо возможностей отклонений. В этом смысле решения \concept{систем ограничений ранга 1} являются доказательствами правильного выполнения программы.

\begin{example}[R1CS для 3-факторизации]\label{ex:3-factorization-r1cs} Чтобы дать лучшую интуицию о \concept{системах ограничений ранга 1}, рассмотрим снова язык $L_{3.fac\_zk}$ из \examplename{} \ref{ex:L-3fac-zk}. Как мы видели, $L_{3.fac\_zk}$ состоит из слов $(<I_1>;<W_1,W_2,W_3>)$ над алфавитом $\F_{13}$, так что $I_1 = W_1\cdot W_2\cdot W_3$. Мы покажем, как переписать проблему как \concept{систему ограничений ранга 1}.

Поскольку R1CS — это системы квадратичных уравнений, выражения вроде $W_1\cdot W_2\cdot W_3$, которые содержат произведения более чем двух множителей (и поэтому не являются квадратичными), должны быть переписаны в процессе, часто называемом \term{выравниванием} (flattening). Чтобы выровнять определяющее уравнение $I_1 = W_1\cdot W_2\cdot W_3$ языка $L_{3.fac\_zk}$, мы вводим новую переменную $W_4$, которая захватывает два из трёх умножений в $W_1\cdot W_2\cdot W_3$. Мы получаем следующие два ограничения
\begin{align*}
W_1 \cdot W_2 & = W_4 & \text{ограничение } 1 \\
W_4 \cdot W_3 & = I_1 & \text{ограничение } 2
\end{align*}
Учитывая некоторый экземпляр $I_1$, любое решение $(W_1,W_2,W_3,W_4)$ этой системы уравнений предоставляет решение исходного уравнения $I_1 = W_1\cdot W_2\cdot W_3$ и наоборот. Оба уравнения поэтому эквивалентны в том смысле, что решения находятся во взаимно однозначном соответствии.

Глядя на оба уравнения из этой системы ограничений, мы видим, как каждое ограничение принуждает шаг в вычислении. Ограничение $1$ принуждает любое вычисление умножить свидетелей $W_1$ и $W_2$; в противном случае было бы невозможно сначала вычислить свидетеля $W_4$, который нужен для решения ограничения $2$. Свидетель $W_4$ поэтому выражает ограничение промежуточного вычислительного состояния.

На этом этапе можно спросить, почему уравнение $1$ ограничивает систему вычислением $W_1\cdot W_2$ первым. Для вычисления $W_1\cdot W_2 \cdot W_3$ вычисление $W_2\cdot W_3$ или $W_1\cdot W_3$ в начале, а затем умножение результата на оставшийся множитель даёт точно такой же результат. Причина чисто вопрос выбора. Например, следующая R1CS определяла бы точно тот же язык:
\begin{align*}
W_2 \cdot W_3 & = W_4 & \text{ограничение } 1 \\
W_4 \cdot W_1 & = I_1 & \text{ограничение } 2 
\end{align*}
Следовательно, R1CS вообще говоря \term{не являются уникальными} описаниями любой данной ситуации: множество различных R1CS способны описать одну и ту же проблему.

Чтобы увидеть, что два квадратичных уравнения квалифицируются как \concept{система ограничений ранга 1}, выберем параметры $n=1$, $m=4$ и $k=2$, а также следующие значения:
$$
\begin{array}{llllll}
a_0^1 = 0 & a_1^1= 0 & a_2^1= 1 & a_3^1 = 0 & a_4^1= 0  & a_5^1= 0 \\ 
a_0^2 = 0 & a_1^2= 0 & a_2^2= 0 & a_3^2 = 0 & a_4^2= 0  & a_5^2= 1 \\ 
\\
b_0^1 = 0 & b_1^1= 0 & b_2^1= 0 & b_3^1 = 1 & b_4^1= 0  & b_5^1= 0 \\ 
b_0^2 = 0 & b_1^2= 0 & b_2^2= 0 & b_3^2 = 0 & b_4^2= 1  & b_5^2= 0 \\ 
\\
c_0^1 = 0 & c_1^1= 0 & c_2^1= 0 & c_3^1 = 0 & c_4^1= 0  & c_5^1= 1 \\ 
c_0^2 = 0 & c_1^2= 1 & c_2^2= 0 & c_3^2 = 0 & c_4^2= 0  & c_5^2= 0 
\end{array} 
$$
С этим выбором \concept{система ограничений ранга 1} нашей проблемы $3$-факторизации может быть записана в наиболее общем виде следующим образом:
\begin{align*}
\scriptstyle
\left(a^1_0 + a_1^1 I_1 + a_2^1 W_1 + a_3^1 W_2 + a_4^1 W_3 + a_5^1 W_4\right)\cdot
\left(b^1_0 + b_1^1 I_1 + b_2^1 W_1 + b_3^1 W_2 + b_4^1 W_3 + b_5^1 W_4\right) &=
\scriptstyle
\left(c^1_0 + c_1^1 I_1 + c_2^1 W_1 + c_3^1 W_2 + c_4^1 W_3 + c_5^1 W_4\right)\\
\scriptstyle
\left(a^2_0 + a_1^2 I_1 + a_2^2 W_1 + a_3^2 W_2 + a_4^2 W_3 + a_5^2 W_4\right)\cdot
\left(b^2_0 + b_1^2 I_1 + b_2^2 W_1 + b_3^2 W_2 + b_4^2 W_3 + b_5^2 W_4\right) &=
\scriptstyle
\left(c^2_0 + c_1^2 I_1 + c_2^2 W_1 + c_3^2 W_2 + c_4^2 W_3 + c_5^2 W_4\right)
\end{align*}
\end{example}
\begin{example}[R1CS для точек кривой \curvename{Tiny-jubjub}]
\label{ex:TJJ-r1cs}
 Рассмотрим языки $L_{tiny.jj.1}$ из \examplename{} \ref{ex:tiny-jubjub-instance-witness}, которые состоят из слов $<I_1,I_2>$ над алфавитом $\F_{13}$, так что $3\cdot I_1^2 + I_2^2 = 1 + 8\cdot I_1^2\cdot I_2^2$. 

Мы выводим \concept{систему ограничений ранга 1} таким образом, что её решения находятся во взаимно однозначном соответствии со словами в $L_{tiny.jj.1}$. Чтобы достичь этого, мы сначала переписываем определяющее уравнение:
\begin{align*}
3\cdot I_1^2 + I_2^2  & = 1 + 8\cdot I_1^2\cdot I_2^2 & \Leftrightarrow \\
 0 & = 1 + 8\cdot I_1^2\cdot I_2^2 - 3\cdot I_1^2 - I_2^2  & \Leftrightarrow \\
 0 & = 1 + 8\cdot I_1^2\cdot I_2^2 + 10\cdot I_1^2 +12\cdot I_2^2
\end{align*}
Поскольку R1CS — это системы квадратичных уравнений, мы должны переформулировать это выражение в систему квадратичных уравнений. Для этого мы должны ввести новые переменные, ограничивающие промежуточные шаги вычисления, и решить, должны ли эти переменные быть переменными экземпляра или свидетеля. Мы решаем объявить все новые переменные переменными свидетеля, и получаем следующие ограничения:
\begin{align*}
I_1 \cdot I_1 & = W_1 & \text{ограничение } 1\\
I_2 \cdot I_2 & = W_2 & \text{ограничение } 2\\
(8 \cdot W_1) \cdot W_2 & = W_3 & \text{ограничение } 3\\
(12\cdot W_2 + W_3 + 10\cdot W_1 + 1)\cdot 1 & = 0 & \text{ограничение } 4
\end{align*}
Чтобы увидеть, что эти четыре квадратичных уравнения квалифицируются как \concept{система ограничений ранга 1} согласно \defname{} \ref{R1CS}, выберем параметры $n=2$, $m=3$, $k=4$ и следующие значения:
$$
\begin{array}{llllll}
a_0^1 = 0 & a_1^1= 1 & a_2^1= 0 & a_3^1 = 0 & a_4^1= 0  & a_5^1= 0 \\ 
a_0^2 = 0 & a_1^2= 0 & a_2^2= 1 & a_3^2 = 0 & a_4^2= 0  & a_5^2= 0 \\ 
a_0^3 = 0 & a_1^3= 0 & a_2^3= 0 & a_3^3 = 8 & a_4^3= 0  & a_5^3= 0 \\ 
a_0^4 = 1 & a_1^4= 0 & a_2^4= 0 & a_3^4 = 10 & a_4^4= 12  & a_5^4= 1 \\ 
\\
b_0^1 = 0 & b_1^1= 1 & b_2^1= 0 & b_3^1 = 0 & b_4^1= 0  & b_5^1= 0 \\ 
b_0^2 = 0 & b_1^2= 0 & b_2^2= 1 & b_3^2 = 0 & b_4^2= 0  & b_5^2= 0 \\ 
b_0^3 = 0 & b_1^3= 0 & b_2^3= 0 & b_3^3 = 0 & b_4^3= 1  & b_5^3= 0 \\ 
b_0^4 = 1 & b_1^4= 0 & b_2^4= 0 & b_3^4 = 0 & b_4^4= 0  & b_5^4= 0 \\ 
\\
c_0^1 = 0 & c_1^1= 0 & c_2^1= 0 & c_3^1 = 1 & c_4^1= 0  & c_5^1= 0 \\ 
c_0^2 = 0 & c_1^2= 0 & c_2^2= 0 & c_3^2 = 0 & c_4^2= 1  & c_5^2= 0 \\
c_0^3 = 0 & c_1^3= 0 & c_2^3= 0 & c_3^3 = 0 & c_4^3= 0  & c_5^3= 1 \\ 
c_0^4 = 0 & c_1^4= 0 & c_2^4= 0 & c_3^4 = 0 & c_4^4= 0  & c_5^4= 0
\end{array} 
$$
С этим выбором \concept{система ограничений ранга 1} нашей проблемы точек кривой \curvename{Tiny-jubjub} может быть записана в наиболее общем виде следующим образом:\tbdsm{font too small}
\begin{align*}
\scriptstyle
\left(a_0^1 + a_1^1 I_1 + a_2^1 I_2 + a_3^1 W_1 + a_4^1 W_2 + a_5^1 W_3\right)\cdot
\left(b_0^1 + b_1^1 I_1 + b_2^1 I_2 + b_3^1 W_1 + b_4^1 W_2 + b_5^1 W_3\right) &=
\scriptstyle
\left(c_0^1 + c_1^1 I_1 + c_2^1 I_2 + c_3^1 W_1 + c_4^1 W_2 + c_5^1 W_3\right)\\
\scriptstyle
\left(a_0^2 + a_1^2 I_1 + a_2^2 I_2 + a_3^2 W_1 + a_4^2 W_2 + a_5^2 W_3\right)\cdot
\left(b_0^2 + b_1^2 I_1 + b_2^2 I_2 + b_3^2 W_1 + b_4^2 W_2 + b_5^2 W_3\right) &=
\scriptstyle
\left(c_0^2 + c_1^2 I_1 + c_2^2 I_2 + c_3^2 W_1 + c_4^2 W_2 + c_5^2 W_3\right)\\\scriptstyle
\left(a_0^3 + a_1^3 I_1 + a_2^3 I_2 + a_3^3 W_1 + a_4^3 W_2 + a_5^3 W_3\right)\cdot
\left(b_0^3 + b_1^3 I_1 + b_2^3 I_2 + b_3^3 W_1 + b_4^3 W_2 + b_5^3 W_3\right) &=
\scriptstyle
\left(c_0^3 + c_1^3 I_1 + c_2^3 I_2 + c_3^3 W_1 + c_4^3 W_2 + c_5^3 W_3\right)\\\scriptstyle
\left(a_0^4 + a_1^4 I_1 + a_2^4 I_2 + a_3^4 W_1 + a_4^4 W_2 + a_5^4 W_3\right)\cdot
\left(b_0^4 + b_1^4 I_1 + b_2^4 I_2 + b_3^4 W_1 + b_4^4 W_2 + b_5^4 W_3\right) &=
\scriptstyle
\left(c_0^4 + c_1^4 I_1 + c_2^4 I_2 + c_3^4 W_1 + c_4^4 W_2 + c_5^4 W_3\right)\\
\end{align*}

Решения этой системы ограничений находятся во взаимно однозначном соответствии со словами в $L_{tiny.jj.1}$: если $(<I_1,I_2>; <W_1, W_2, W_3>)$ является решением, то $<I_1,I_2>$ является словом в $L_{tiny.jj.1}$, поскольку определяющая R1CS подразумевает, что $I_1$ и $I_2$ удовлетворяют уравнению скрученного Эдвардса кривой \curvename{Tiny-jubjub}. С другой стороны, если $<I_1,I_2>$ является словом в $L_{tiny.jj.1}$, то $(<I_1,I_2>; <I_1^2, I_2^2, 8\cdot I_1^2\cdot I_2^2>)$ является решением нашей R1CS.
\end{example}
\begin{exercise} Рассмотрим язык $L_{add}$ из упражнения \ref{ex:TJJ-addition-lang}. Определите R1CS таким образом, что слова в $L_{add}$ находятся во взаимно однозначном соответствии с решениями этой R1CS.
\end{exercise} 
\subsubsection{Выполнимость R1CS}
\label{r1cs-constructive-proofs}
Чтобы понять, как \concept{системы ограничений ранга 1} определяют формальные языки, заметим, что каждая R1CS над полем $\F$ определяет функцию решения над алфавитом $\Sigma_I \times \Sigma_W = \F \times \F$ следующим образом:
\begin{equation}
R_{R1CS} : (\F)^* \times (\F)^* \to \{true, false\}\;;\;
(I;W) \mapsto
\begin{cases}
true & (I;W) \text{ удовлетворяет R1CS}\\
false & \text{ иначе}
\end{cases}
\end{equation}
Каждая R1CS поэтому определяет формальный язык. Грамматика этого языка закодирована в ограничениях, слова — это решения уравнений, а \term{утверждение} (statement) — это утверждение о знании ``Учитывая экземпляр $I$, существует свидетель $W$, такой что $(I;W)$ является решением \concept{системы ограничений ранга 1}''. Конструктивное доказательство этого утверждения поэтому эквивалентно присваиванию элемента поля каждой переменной свидетеля, что верифицируется всякий раз, когда множество всех переменных экземпляра и свидетеля решает R1CS.

\begin{remark}[Выполнимость R1CS]\label{r1cs-satisfiability}  \sme{SB: delete this or move to a footnote}Следует отметить, что в нашем определении каждая R1CS определяет свой собственный язык. Однако в более теоретических подходах часто рассматривается другой язык, обычно называемый \term{выполнимостью R1CS} (R1CS satisfiability), который полезен, когда речь идёт о более абстрактных проблемах, таких как выразительность или вычислительная сложность класса \term{всех} R1CS. С нашей точки зрения, язык выполнимости R1CS получается объединением всех языков R1CS, которые есть в нашем определении. Будем более точны. Пусть алфавит $\Sigma=\F$ — поле. Тогда язык $L_{R1CS\_SAT(\F)}$ определяется следующим образом:
$$
L_{R1CS\_SAT(\F)} = \{(i;w)\in \Sigma^*\times \Sigma^*\;|\; \text{существует R1CS } R \text{ такая, что } R(i;w)=true  \}
$$
\end{remark}
\begin{example}[3-Factorization]
\label{ex:3-fac-R1CS-constr-proof}
Рассмотрим язык $L_{3.fac\_zk}$ из \examplename{} \ref{ex:L-3fac-zk} и R1CS, определённую в \examplename{} \ref{ex:3-factorization-r1cs}. Как мы видели в \ref{ex:3-factorization-r1cs}, решения R1CS находятся во взаимно однозначном соответствии с решениями функции решения $L_{3.fac\_zk}$. Оба языка поэтому эквивалентны в том смысле, что существует взаимно однозначное соответствие между словами в обоих языках.

Чтобы дать интуицию о том, как выглядят конструктивные R1CS-based доказательства в $L_{3.fac\_zk}$, рассмотрим экземпляр $I_1= 11$. Чтобы доказать утверждение ``Существует свидетель $W$, такой что $(I_1;W)$ является словом в $L_{3.fac\_zk}$'' конструктивно, доказательство должно предоставить решение R1CS из \examplename{} \ref{ex:3-factorization-r1cs}, то есть присваивание всем переменным свидетеля $W_1$, $W_2$, $W_3$ и $W_4$. Поскольку алфавитом является $\F_{13}$, пример присваивания даётся
$W=<2,3,4,6>$ поскольку $(I_1;W)$ удовлетворяет R1CS:
\begin{align*}
W_1 \cdot W_2 &= W_4 & \text{\# } 2\cdot 3 = 6\\
W_4 \cdot W_3 &= I_1 & \text{\# } 6\cdot 4 = 11
\end{align*}
Правильное конструктивное доказательство поэтому даётся $\pi=<2,3,4,6>$. Конечно, $\pi$ — не единственное возможное доказательство этого утверждения. Поскольку разложение на множители не является уникальным в поле вообще говоря, другое конструктивное доказательство даётся $\pi'=<3,5,12,2>$. 

\begin{align*}
W_1 \cdot W_2 &= W_4 & \text{\# } 3\cdot 5 = 2\\
W_4 \cdot W_3 &= I_1 & \text{\# } 12\cdot 2 = 11
\end{align*}

\end{example}

\begin{example}[The \curvename{Tiny-jubjub} curve]\label{ex:tiny-jubjub-r1cs} Рассмотрим язык $L_{tiny.jj.1}$ из \examplename{} \ref{ex:tiny-jubjub-instance-witness}, и ассоциированную R1CS из \examplename{} \ref{ex:TJJ-r1cs}. Чтобы увидеть, как выглядят конструктивные доказательства в $L_{tiny.jj.1}$ с использованием R1CS из \examplename{} \ref{ex:TJJ-r1cs}, рассмотрим экземпляр $<I_1,I_2>= <11,6>$. Чтобы доказать утверждение ``Существует свидетель $W$, такой что $(<I_1,I_2>;W)$ является словом в $L_{tiny.jj.1}$'' конструктивно, доказательство должно предоставить решение R1CS \ref{ex:TJJ-r1cs}, что является присваиванием всем переменным свидетеля $W_1$, $W_2$ и $W_3$. Поскольку алфавитом является $\F_{13}$, пример присваивания даётся
$W=<4,10,8>$ поскольку $(<I_1,I_2>;W)$ удовлетворяет R1CS:
\begin{align*}
I_1 \cdot I_1 & = W_1 & 11\cdot 11 = 4\\
I_2 \cdot I_2 & = W_2 & 6 \cdot 6 = 10 \\
(8 \cdot W_1) \cdot W_2 & = W_3 & (8\cdot 4)\cdot 10 = 8\\
(12\cdot W_2 + W_3 + 10\cdot W_1 + 1)\cdot 1 & = 0 & 12\cdot 10 + 8 + 10\cdot 4 + 1 = 0
\end{align*}
Правильное конструктивное доказательство поэтому даётся $\pi=<4,10,8>$, что показывает, что экземпляр $<11,6>$ является точкой на кривой \curvename{Tiny-jubjub}. 
\end{example}
\subsubsection{Модульность}
\label{sec:R1CS_modularity} Как мы обсуждали в \ref{modularity}, часто полезно конструировать сложные утверждения и представляющие их языки из простых. \concept{Системы ограничений ранга 1} особенно полезны для этого, поскольку пересечение двух R1CS над одним и тем же алфавитом даёт новую R1CS над тем же алфавитом.

Будем более точны. Пусть $S_1$ и $S_2$ — две R1CS над $\F$. Новая R1CS $S_3$ получается пересечением $S_1$ и $S_2$, то есть $S_3 = S_1\cap S_2$. В этом контексте пересечение означает, что должны быть удовлетворены как уравнения $S_1$ \term{и} уравнения $S_2$, чтобы предоставить решение для системы $S_3$.

Как следствие, разработчики способны конструировать сложные R1CS из простых. Эта модульность предоставляет теоретический фундамент для многих компиляторов R1CS, как мы увидим в \chaptname{} \ref{chap:circuit-compilers}.

\subsection{Алгебраические схемы}
\label{sec:circuits} Как мы видели в предыдущих параграфах, \concept{системы ограничений ранга 1} являются квадратичными уравнениями, так что решения являются доказательствами знания существования слов в ассоциированных языках. С точки зрения доказывающего, поэтому важно эффективно решать эти уравнения.

Однако, в отличие от систем линейных уравнений, неизвестны общие методы, эффективно решающие системы квадратичных уравнений. \concept{Системы ограничений ранга 1} поэтому непрактичны с точки зрения доказывающего, и необходима вспомогательная информация, чтобы помочь эффективно вычислять решения.

Методы, вычисляющие решения R1CS, иногда называются \term{функциями генерации свидетелей} (witness generator functions). Чтобы привести распространённый пример, мы вводим другой класс функций решения, называемых \term{алгебраическими схемами} (algebraic circuits). Как мы увидим, каждая алгебраическая схема определяет ассоциированную R1CS и также предоставляет эффективный способ вычисления решений для этой R1CS. Этот метод введён, например, в \cite{sasson-2013}.

Можно показать, что любое ограниченное по пространству и времени вычисление может быть выражено как алгебраическая схема. Преобразование высокоуровневых компьютерных программ в эти схемы — процесс, часто называемый \term{выравниванием} (flattening). Мы рассмотрим эти преобразования в \chaptname{} \ref{chap:circuit-compilers}.

В этом разделе мы введём нашу модель для алгебраических схем и рассмотрим концепцию выполнения схемы (circuit execution) и корректных присваиваний (valid assignments). После этого мы покажем, как выводить \concept{системы ограничений ранга 1} из схем и как схемы полезны для эффективного вычисления решений ассоциированных R1CS.
\subsubsection{Представление алгебраической схемы} Чтобы увидеть, что такое алгебраические схемы, пусть $\F$ — поле. Алгебраическая схема тогда — это ориентированный ациклический мультиграф (directed acyclic multi-graph), вычисляющий полиномиальную функцию над $\F$. Узлы только с исходящими рёбрами (источники, source nodes) представляют переменные и константы функции, а узлы только с входящими рёбрами (стоки, sink nodes) представляют результат функции. Все остальные узлы имеют ровно два входящих ребра и представляют операции поля \term{сложение} (addition), а также \term{умножение} (multiplication). Рёбра графа ориентированы и представляют поток вычисления вдоль узлов.

Будем более точны. В этой книге мы называем ориентированный ациклический мультиграф $C(\F)$ \term{алгебраической схемой} (algebraic circuit) над $\F$, если выполняются следующие условия:

\begin{itemize}
\label{def:algebraic-circuit}
\item Множество рёбер имеет полный порядок.
\item Каждый исходный узел имеет метку, представляющую либо переменную, либо константу из поля $\F$.
\item Каждый стоковый узел имеет ровно одно входящее ребро и метку, представляющую либо переменную, либо константу из поля $\F$.
\item Каждый узел, который не является ни исходным, ни стоковым, имеет ровно два входящих ребра и метку из множества $\{+,*\}$, представляющую либо сложение, либо умножение в $\F$.
\item Все исходящие рёбра из узла имеют одну и ту же метку.
\item Исходящие рёбра из узла с меткой, представляющей переменную, имеют метку.
\item Исходящие рёбра из узла с меткой, представляющей умножение, имеют метку, если существует хотя бы одно помеченное ребро в обоих входных путях.
\item Все входящие рёбра в стоковые узлы имеют метку.
\item Если ребро имеет две метки $S_i$ и $S_j$, оно получает новую метку $S_i = S_j$.
\item Никакое другое ребро не имеет метки.
\item Входящие рёбра в помеченные стоковые узлы, где метка является константой $c\in\F$, помечены той же константой. Каждая другая метка ребра берётся из множества $\{W,I\}$ и индексируется совместно с порядком множества рёбер. 
\end{itemize} 

Следует отметить, что детали в определениях алгебраических схем варьируются между различными источниками. Мы используем это определение, поскольку оно концептуально простое и хорошо подходит для вычислений от руки.

Чтобы получить лучшую интуицию нашего определения, пусть $C(\F)$ — алгебраическая схема. Исходные узлы — это входы схемы и либо представляют переменные, либо константы. Аналогичным образом, стоковые узлы представляют точки завершения схемы и либо являются выходными переменными, либо константами. Константные стоковые узлы принуждают вычислительные выходы принимать определённые значения.

Узлы, которые не являются ни исходными, ни стоковыми, называются \term{арифметическими вентилями} (arithmetic gates). Арифметические вентили, помеченные меткой ``$+$'', называются \term{вентилями сложения} (addition-gates), а арифметические вентили, помеченные меткой ``$\cdot$'', называются \term{вентилями умножения} (multiplication-gates). Каждый арифметический вентиль имеет ровно два входа, представленных двумя входящими рёбрами.

Поскольку множество рёбер упорядочено, мы можем записать его как строку $<E_1,E_2,\ldots, E_n>$ для некоторого $n\in \N$, и мы используем эти индексы для индексации меток рёбер тоже. Метки рёбер поэтому либо константы, либо символы вроде $I_j$, $W_j$, где $j$ — индекс, совместимый с порядком рёбер. Метки $I_j$ представляют переменные экземпляра, метки $W_j$ — переменные свидетеля. Метки на исходящих рёбрах входных переменных связывают ассоциированную переменную с этим ребром.

\begin{notation}
При синтезе алгебраических схем присваивание меток экземпляра $I_j$ или свидетеля $W_j$ соответствующим рёбрам часто является финальным шагом. Поэтому удобно не различать эти два типа рёбер на предыдущих шагах. Чтобы учесть это, мы часто просто пишем $S_j$ для метки ребра, указывая, что свойство экземпляра/свидетеля метки не определено, и она может представлять как метку экземпляра, так и метку свидетеля. 
\end{notation}
\begin{example}[Обобщённая факторизация SNARK]
\label{ex:3-fac-zk-circuit} Чтобы привести простой пример алгебраической схемы, рассмотрим снова нашу проблему $3$-факторизации из \examplename{} \ref{ex:L-3fac-zk}. Чтобы выразить проблему в модели алгебраической схемы, рассмотрим следующую функцию 
\[
f_{3.fac}:\mathbb{F}_{13}\times\mathbb{F}_{13}\times\mathbb{F}_{13}\to\mathbb{F}_{13};(x_{1},x_{2},x_{3})\mapsto x_{1}\cdot x_{2}\cdot x_{3}
\]
Используя эту функцию, мы можем описать проблему $3$-факторизации с нулевым разглашением из \ref{ex:L-3fac-zk} следующим образом: Учитывая экземпляр $I_1\in \F_{13}$, корректный свидетель — это прообраз $f_{3.fac}$ в точке $I_1$, то есть корректный свидетель состоит из трёх значений $W_1$, $W_2$ и $W_3$ из $\F_{13}$, таких что $f_{3.fac}(W_1,W_2,W_3)=I_1$. 

Чтобы увидеть, как эта функция может быть преобразована в алгебраическую схему над $\F_{13}$, общим первым шагом является введение скобок в определение функции, а затем запись операций как бинарных операторов, чтобы подчеркнуть, как именно каждая операция поля действует на свои два входа. В силу законов ассоциативности в поле у нас есть несколько вариантов. Мы выбираем
\begin{align*}
f_{3.fac}(x_1,x_2,x_3) & = x_1\cdot x_2 \cdot x_3  & \text{\# bracket choice} \\
                       & = (x_1\cdot x_2 ) \cdot x_3  & \text{\# operator notation} \\
                       & = MUL(MUL(x_1,x_2),x_3)
\end{align*}
Используя это выражение, мы можем записать ассоциированную алгебраическую схему, сначала связывая переменные с метками рёбер $W_1=x_1$, $W_2=x_2$ и $W_3=x_3$, а также $I_1=f_{3.fac}(x_1,x_2,x_3)$, учитывая различие между входами свидетеля и экземпляра. Затем мы переписываем операторное представление $f_{3.fac}$ в узлы схемы и получаем следующее: 
\begin{center}
% https://www.graphviz.org/pdf/dotguide.pdf
\digraph[scale=0.5]{G1}{
	forcelabels=true;
	//center=true;
	splines=ortho;
	nodesep= 2.0;
	n1 -> n3 [xlabel="W_2  "];
	n2 -> n3 [xlabel="  W_1"];
	n3 -> n5 [label="W_4"];
	n4 -> n5 [label=" W_3"];
	n5 -> n6 [label=" I_1"];
	n1 [shape=box, label="x_2"];
	n2 [shape=box, label="x_1"];
	n3 [label="*"];
	n4 [shape=box, label="x_3"];
	n5 [label="*"];
	n6 [shape=box, label="f_(3.fac_zk)"];
}
\end{center}
%\[
%\xymatrix{x_1\ar^{W_1}[dr] &  & x_2\ar_{W_2}[dl]\\
% & \cdot \ar^{W_4}[drr] &   & & x_3\ar_{W_3}[dl]\\
%  &  &  & \cdot\ar_{I_1}[d]\\
%  &  &  & f(x_1,x_2,x_3)
%}
%\]


В этом случае ориентированный ациклический мультиграф является бинарным деревом с тремя листьями (исходными узлами), помеченными $x_1$, $x_2$ и $x_3$, одним корнем (единственным стоковым узлом), помеченным $f_{3.fac}(x_1,x_2,x_3)$, и двумя внутренними узлами, которые помечены как вентили умножения. 

Порядок, который мы используем для маркировки рёбер, выбран так, чтобы маркировка рёбер была согласована с выбором $W_4$, как определено в определении \ref{def:algebraic-circuit}. Этот порядок может быть получен алгоритмом обхода в глубину справа налево.
\end{example}
\begin{example}
\label{ex:TJJ-circuit_1} Чтобы привести более реалистичный пример алгебраической схемы, снова посмотрим на определяющее уравнение кривой \curvename{Tiny-jubjub} \ref{TJJ13}. Пара элементов поля 
$(x,y)\in \F_{13}^2$ является точкой кривой, точно если выполняется следующее уравнение:
$$
3\cdot x^2 + y^2 = 1+ 8\cdot x^2\cdot y^2
$$ 
Чтобы понять, как можно преобразовать это тождество в алгебраическую схему, мы сначала переписываем это уравнение, сдвигая все члены вправо. Мы получаем следующее:
\begin{align*}
3\cdot x^2 + y^2 & = 1+ 8\cdot x^2\cdot y^2 & \Leftrightarrow\\
0 & = 1+ 8\cdot x^2\cdot y^2 - 3\cdot x^2 - y^2 & \Leftrightarrow\\
0 & = 1+ 8\cdot x^2\cdot y^2 + 10\cdot x^2 +12\cdot y^2
\end{align*}
Затем мы используем это выражение, чтобы определить функцию так, что все точки кривой \curvename{Tiny-jubjub} характеризуются как прообразы функции в точке $0$.
$$
f_{tiny-jj}: \F_{13}\times \F_{13}\to \F_{13}\; ; \;
(x,y)\mapsto 1+ 8\cdot x^2\cdot y^2 + 10\cdot x^2 +12\cdot y^2
$$
Каждая пара элементов поля $(x,y)\in \F_{13}^2$ с $f_{tiny-jj}(x,y)=0$ является точкой на кривой \curvename{Tiny-jubjub}, и других точек кривой не существует. Прообраз $f_{tiny-jj}^{-1}(0)$ поэтому является полным описанием кривой \curvename{Tiny-jubjub}.

Мы можем преобразовать эту функцию в алгебраическую схему над $\F_{13}$. Мы сначала вводим скобки в потенциально неоднозначные выражения, а затем переписываем функцию в терминах бинарных операторов. Мы получаем следующее:
\begin{align*}
\label{ex:tiny_circuit_brackets}
f_{tiny-jj}(x,y) & = 1+ 8\cdot x^2\cdot y^2 + 10\cdot x^2 +12 y^2  & \Leftrightarrow\\
 & = ((8\cdot ((x\cdot x)\cdot (y\cdot y))) + (1+ 10\cdot (x\cdot x))) + (12\cdot(y\cdot y))  & \Leftrightarrow\\
  & = \scriptstyle ADD(ADD(MUL(8,MUL(MUL(x,x),MUL(y,y))), ADD(1,MUL(10,MUL(x,x)))),MUL(12,MUL(y,y)))
\end{align*}\tbdsm{this needs to be numbered}
Поскольку мы ещё не решили, какая часть вычисления должна быть экземпляром, а какая — свидетелем, мы используем неопределённый символ $S$ для представления меток рёбер. Связывая все переменные с метками рёбер $S_1=x$, $S_2=y$ и $S_6=f_{tiny-jj}$, мы получаем следующую схему, представляющую функцию $f_{tiny-jj}$, индуктивно заменяя бинарные операторы их ассоциированными арифметическими вентилями:
\begin{center}
% https://www.graphviz.org/pdf/dotguide.pdf
\digraph[scale=0.4]{G2}{
	forcelabels=true;
	//center=true;
	splines=ortho;
	nodesep= 2.0;
	n1 -> n4 [xlabel="S_1" /*, color=lightgray */];
        n1 -> n4 /*[ color=lightgray ]*/
	n2 -> n5 [xlabel="S_2" /*, color=lightgray */];
	n2 -> n5 /*[ color=lightgray ]*/ ;
	n3 -> n8 /*[ color=lightgray ]*/ ;
	n4 -> n8 [taillabel="S_3", labeldistance="2" /*, color=lightgray */];
	n4 -> n10 [taillabel="S_3", labeldistance="4" /*, color=lightgray */];
	n5 -> n10 [xlabel="S_4" /*, color=lightgray */];
	n5 -> n13 [xlabel="S_4" /*, color=lightgray */];
	n6 -> n13 /*[ color=lightgray ]*/ ;
	n7 -> n11 /*[ color=lightgray ]*/ ;
	n8 -> n11 /*[ color=lightgray ]*/ ;
	n9 -> n12 /*[ color=lightgray ]*/ ;
	n10 -> n12 [taillabel="S_5", labeldistance="4" /*, color=lightgray */];
	n11 -> n14 /*[ color=lightgray ]*/ ;
	n12 -> n14 /*[ color=lightgray ]*/ ;	
	n13 -> n15 /*[ color=lightgray ]*/ ;
	n14 -> n15 /*[ color=lightgray ]*/ ;
	n15 -> n16 [label="  S_6", labeldistance="2" /*, color=lightgray */];
	n1 [shape=box, label="x" /*, color=lightgray */];
	n2 [shape=box, label="y" /*, color=lightgray */];
	n3 [shape=box, label="10" /*, color=lightgray */];
	n4 [label="*" /*, color=lightgray */];
	n5 [label="*" /*, color=lightgray */];
	n6 [shape=box, label="12" /*, color=lightgray */];
	n7 [shape=box, label="1" /*, color=lightgray */];
	n8 [label="*" /*, color=lightgray */];
	n9 [shape=box, label="8" /*, color=lightgray */];
	n10 [label="*" /*, color=lightgray */];
	n11 [label="+" /*, color=lightgray */];
	n12 [label="*" /*, color=lightgray */];	
	n13 [label="*" /*, color=lightgray */];
	n14 [label="+" /*, color=lightgray */];
	n15 [label="+" /*, color=lightgray */];
	n16 [shape=box, label="f_tiny-jj" /*, color=lightgray */];		
}
\end{center}
Эта схема не является графом, а мультиграфом, поскольку между некоторыми узлами существует более одного ребра. 

В процессе проектирования схем из функций следует отметить, что представления схем в общем случае не являются уникальными. В случае функции $f_{tiny-jj}$ форма схемы зависит от нашего выбора скобок выше.%in \ref{ex:tiny_circuit_brackets}. 
Альтернативный дизайн, например, даётся следующей схемой, которая возникает, когда скобочное выражение $8\cdot ( (x\cdot x) \cdot (y\cdot y) )$ заменяется выражением $(x\cdot x) \cdot ( 8 \cdot (y\cdot y) )$. 
\begin{center}
\digraph[scale=0.4]{G2A}{
	forcelabels=true;
	center=true;
	splines=ortho;
	nodesep= 2.0;
	n1 -> n4 [label="S_1" /* comment*/ labeldistance="4" /*, color=lightgray */];
        n1 -> n4 /*[ color=lightgray ]*/
	n2 -> n6 [label="S_2" /*, color=lightgray */];
	n2 -> n6 /*[ color=lightgray ]*/ ;
	n3 -> n9 /*[ color=lightgray ]*/ ;
	n4 -> n9 [taillabel="S_3", labeldistance="2" /*, color=lightgray */];
	n4 -> n12 [taillabel="S_3", labeldistance="4" /*, color=lightgray */];
	n5 -> n10 /*[ color=lightgray ]*/ ;
	n6 -> n10 [headlabel="S_4", labeldistance="4" /*, color=lightgray */];
	n6 -> n13 [taillabel="S_4", labeldistance="4" /*, color=lightgray */];
	n7 -> n13 /*[ color=lightgray ]*/ ;
	n8 -> n11 /*[ color=lightgray ]*/ ;
	n9 -> n11 /*[ color=lightgray ]*/ ;
	n10 -> n12 /*[ color=lightgray ]*/ ; 
	n11 -> n14 /*[ color=lightgray ]*/ ;
	n12 -> n14 [xlabel="S_5  "  /*, color=lightgray */];	
	n13 -> n15 /*[ color=lightgray ]*/ ;
	n14 -> n15 /*[ color=lightgray ]*/ ;
	n15 -> n16 [label="  S_6", labeldistance="2" /*, color=lightgray */];
	n1 [shape=box, label="x" /*, color=lightgray */];
	n2 [shape=box, label="y" /*, color=lightgray */];
	n3 [shape=box, label="10" /*, color=lightgray */];
	n4 [label="*" /*, color=lightgray */];
	n5 [shape=box, label="8" /*, color=lightgray */];
	n6 [label="*" /*, color=lightgray */];
	n7 [shape=box, label="12" /*, color=lightgray */];
	n8 [shape=box, label="1" /*, color=lightgray */];
	n9 [label="*" /*, color=lightgray */];
	n10 [label="*" /*, color=lightgray */];
	n11 [label="+" /*, color=lightgray */];	
	n12 [label="*" /*, color=lightgray */];	
	n13 [label="*" /*, color=lightgray */];
	n14 [label="+" /*, color=lightgray */];
	n15 [label="+" /*, color=lightgray */];
	n16 [shape=box, label="f_tiny-jj" /*, color=lightgray */];		
}
\end{center}
Конечно, обе схемы представляют одну и ту же функцию, благодаря законам ассоциативности и коммутативности, которые выполняются в любом поле.

Имея схему, представляющую функцию $f_{tiny-jj}$, мы можем теперь перейти к выводу схемы, ограничивающей произвольные пары $(x,y)$ элементов поля, чтобы они были точками на кривой \curvename{Tiny-jubjub}. Для этого мы должны ограничить выход, чтобы он был равен нулю, то есть мы должны ограничить $S_6=0$. Чтобы указать это в схеме, мы заменяем выходную переменную константой $0$ и связываем соответствующую метку ребра соответственно. Мы получаем следующее:
\begin{center}
\digraph[scale=0.4]{G2AJA}{
	forcelabels=true;
	center=true;
	splines=ortho;
	nodesep= 2.0;
	n1 -> n4 [label="S_1" /* comment*/ labeldistance="4" /*, color=lightgray */];
        n1 -> n4 /*[ color=lightgray ]*/
	n2 -> n6 [label="S_2" /*, color=lightgray */];
	n2 -> n6 /*[ color=lightgray ]*/ ;
	n3 -> n9 /*[ color=lightgray ]*/ ;
	n4 -> n9 [taillabel="S_3", labeldistance="2" /*, color=lightgray */];
	n4 -> n12 [taillabel="S_3", labeldistance="4" /*, color=lightgray */];
	n5 -> n10 /*[ color=lightgray ]*/ ;
	n6 -> n10 [headlabel="S_4", labeldistance="4" /*, color=lightgray */];
	n6 -> n13 [taillabel="S_4", labeldistance="4" /*, color=lightgray */];
	n7 -> n13 /*[ color=lightgray ]*/ ;
	n8 -> n11 /*[ color=lightgray ]*/ ;
	n9 -> n11 /*[ color=lightgray ]*/ ;
	n10 -> n12 /*[ color=lightgray ]*/ ; 
	n11 -> n14 /*[ color=lightgray ]*/ ;
	n12 -> n14 [xlabel="S_5  "  /*, color=lightgray */];	
	n13 -> n15 /*[ color=lightgray ]*/ ;
	n14 -> n15 /*[ color=lightgray ]*/ ;
	n15 -> n16 [label="  S_6=0", labeldistance="2" /*, color=lightgray */];
	n1 [shape=box, label="x" /*, color=lightgray */];
	n2 [shape=box, label="y" /*, color=lightgray */];
	n3 [shape=box, label="10" /*, color=lightgray */];
	n4 [label="*" /*, color=lightgray */];
	n5 [shape=box, label="8" /*, color=lightgray */];
	n6 [label="*" /*, color=lightgray */];
	n7 [shape=box, label="12" /*, color=lightgray */];
	n8 [shape=box, label="1" /*, color=lightgray */];
	n9 [label="*" /*, color=lightgray */];
	n10 [label="*" /*, color=lightgray */];
	n11 [label="+" /*, color=lightgray */];	
	n12 [label="*" /*, color=lightgray */];	
	n13 [label="*" /*, color=lightgray */];
	n14 [label="+" /*, color=lightgray */];
	n15 [label="+" /*, color=lightgray */];
	n16 [shape=box, label="0" /*, color=lightgray */];		
}
\end{center}
Предыдущая схема принуждает входные значения, присвоенные меткам $S_1$ и $S_2$, быть точками на кривой \curvename{Tiny-jubjub}. Однако она не указывает, какие метки считаются экземпляром, а какие — свидетелем. Следующая схема определяет входы как экземпляры, в то время как все остальные метки представляют свидетелей:
\begin{center}
\digraph[scale=0.4]{G2AJA2}{
	forcelabels=true;
	center=true;
	splines=ortho;
	nodesep= 2.0;
	n1 -> n4 [label="I_1" /* comment*/ labeldistance="4" /*, color=lightgray */];
        n1 -> n4 /*[ color=lightgray ]*/
	n2 -> n6 [label="I_2" /*, color=lightgray */];
	n2 -> n6 /*[ color=lightgray ]*/ ;
	n3 -> n9 /*[ color=lightgray ]*/ ;
	n4 -> n9 [taillabel="W_1", labeldistance="2" /*, color=lightgray */];
	n4 -> n12 [taillabel="W_1", labeldistance="4" /*, color=lightgray */];
	n5 -> n10 /*[ color=lightgray ]*/ ;
	n6 -> n10 [headlabel="W_2", labeldistance="4" /*, color=lightgray */];
	n6 -> n13 [taillabel="W_2", labeldistance="4" /*, color=lightgray */];
	n7 -> n13 /*[ color=lightgray ]*/ ;
	n8 -> n11 /*[ color=lightgray ]*/ ;
	n9 -> n11 /*[ color=lightgray ]*/ ;
	n10 -> n12 /*[ color=lightgray ]*/ ; 
	n11 -> n14 /*[ color=lightgray ]*/ ;
	n12 -> n14 [xlabel="W_3  "  /*, color=lightgray */];	
	n13 -> n15 /*[ color=lightgray ]*/ ;
	n14 -> n15 /*[ color=lightgray ]*/ ;
	n15 -> n16 [label="  0", labeldistance="2" /*, color=lightgray */];
	n1 [shape=box, label="x" /*, color=lightgray */];
	n2 [shape=box, label="y" /*, color=lightgray */];
	n3 [shape=box, label="10" /*, color=lightgray */];
	n4 [label="*" /*, color=lightgray */];
	n5 [shape=box, label="8" /*, color=lightgray */];
	n6 [label="*" /*, color=lightgray */];
	n7 [shape=box, label="12" /*, color=lightgray */];
	n8 [shape=box, label="1" /*, color=lightgray */];
	n9 [label="*" /*, color=lightgray */];
	n10 [label="*" /*, color=lightgray */];
	n11 [label="+" /*, color=lightgray */];	
	n12 [label="*" /*, color=lightgray */];	
	n13 [label="*" /*, color=lightgray */];
	n14 [label="+" /*, color=lightgray */];
	n15 [label="+" /*, color=lightgray */];
	n16 [shape=box, label="0" /*, color=lightgray */];		
}
\end{center} 
\end{example}
Можно показать, что любое ограниченное по пространству и времени вычисление может быть преобразовано в алгебраическую схему. Мы называем любой процесс, преобразующий ограниченное вычисление в схему, \term{выравниванием} (flattening).

\subsubsection{Выполнение схемы} Алгебраические схемы — ориентированные ациклические мультиграфы, где узлы представляют переменные, константы или вентили сложения и умножения. В частности, каждая алгебраическая схема с $n$ входными узлами, помеченными символами переменных, и $m$ выходными узлами, помеченными переменными, может рассматриваться как функция, преобразующая входную строку $(x_1,\ldots, x_n)$ из $\F^n$ в выходную строку $(f_1,\ldots,f_m)$ из $\F^m$. Преобразование выполняется отправкой значений, ассоциированных с узлами, вдоль их исходящих рёбер к другим узлам. Если эти узлы являются вентилями, то значения преобразуются согласно метке вентиля, и процесс повторяется вдоль всех рёбер до достижения стокового узла. Мы называем это вычисление \term{выполнением схемы} (circuit execution).

При выполнении схемы возможно не только вычислить выходные значения схемы, но и вывести элементы поля для всех рёбер, и, в частности, для всех меток рёбер в схеме. Результатом является строка $<S_1,S_2,\ldots, S_n>$ элементов поля, ассоциированных со всеми помеченными рёбрами, которую мы называем \term{корректным присваиванием} (valid assignment) схеме. Напротив, любое присваивание $<S'_1,S'_2,\ldots, S'_n>$ элементов поля меткам рёбер, которое не может возникнуть из выполнения схемы, называется \term{некорректным присваиванием} (invalid assignment).

Корректные присваивания могут интерпретироваться как \term{доказательства правильного выполнения схемы}, поскольку они сохраняют запись вычислительного результата, а также промежуточных вычислительных шагов.
\begin{example}[3-factorization]
\label{ex:3-fac-zk-circuit_2}
Рассмотрим проблему $3$-факторизации из \examplename{} \ref{ex:L-3fac-zk} и её представление как алгебраической схемы из \examplename{} \ref{ex:3-fac-zk-circuit}. Мы знаем, что строка меток рёбер даётся $S:=<I_{1};W_{1},W_{2},W_{3}, W_{4}>$. 

Чтобы понять, как выполняется эта схема, рассмотрим переменные $x_1=2$, $x_2=3$, а также $x_3=4$. Следуя всем рёбрам в графе, мы получаем присваивания $W_1=2$, $W_2=3$ и $W_3=4$. Затем присваивания $W_1$ и $W_2$ входят в вентиль умножения, и выход вентиля равен $2\cdot 3 = 6$, что мы присваиваем $W_4$, то есть $W_4=6$. Значения $W_4$ и $W_3$ затем входят во второй вентиль умножения, и выход вентиля равен $6\cdot 4 = 11$, что мы присваиваем $I_1$, то есть $I_1=11$. 

Корректное присваивание схеме $3$-факторизации $C_{3.fac}(\F_{13})$ поэтому даётся следующей строкой элементов поля из $\F_{13}$:

\begin{equation}\label{C3fac}
S_{valid}:=<11;2,3,4,6>
\end{equation}

Мы можем визуализировать это присваивание, присваивая каждое вычисленное значение его ассоциированной метке в схеме следующим образом:
\begin{center}
% https://www.graphviz.org/pdf/dotguide.pdf
\digraph[scale=0.5]{G3}{
	forcelabels=true;
	//center=true;
	splines=ortho;
	nodesep= 2.0;
	n1 -> n3 [xlabel="W_2=3  "];
	n2 -> n3 [xlabel="W_1=2 "];
	n3 -> n5 [label="W_4=6"];
	n4 -> n5 [label=" W_3=4"];
	n5 -> n6 [label=" I_1=11"];
	n1 [shape=box, label="x_2"];
	n2 [shape=box, label="x_1"];
	n3 [label="*"];
	n4 [shape=box, label="x_3"];
	n5 [label="*"];
	n6 [shape=box, label="f_(3.fac_zk)"];
}
\end{center}
Чтобы увидеть, как выглядит некорректное присваивание, рассмотрим присваивание $S_{err}:=<8;2,3,4,7>$. В этом присваивании входные значения такие же, как в предыдущем случае. Ассоциированная схема:
\begin{center}
% https://www.graphviz.org/pdf/dotguide.pdf
\digraph[scale=0.5]{G4}{
	forcelabels=true;
	//center=true;
	splines=ortho;
	nodesep= 2.0;
	n1 -> n3 [xlabel="W_2=3  "];
	n2 -> n3 [xlabel="W_1=2 "];
	n3 -> n5 [label="W_4=7"];
	n4 -> n5 [label=" W_3=4"];
	n5 -> n6 [label=" I_1=8"];
	n1 [shape=box, label="x_2"];
	n2 [shape=box, label="x_1"];
	n3 [label="*"];
	n4 [shape=box, label="x_3"];
	n5 [label="*"];
	n6 [shape=box, label="f_(3.fac_zk)"];
}
\end{center}
Это присваивание является некорректным, поскольку присваивания $I_1$ и $W_4$ не могут быть получены выполнением схемы.
\end{example}
\begin{example}
\label{ex:TJJ-circuit_2} 
 Чтобы вычислить более реалистичное выполнение алгебраической схемы, снова рассмотрим определяющую схему $C_{tiny-jj}(\F_{13})$ из \examplename{} \ref{ex:TJJ-circuit_1}. Мы уже знаем из способа построения этой схемы, что любое корректное присваивание с $S_1=x$, $S_2=y$ и $S_6=0$ гарантирует, что пара $(x,y)$ является точкой на кривой \curvename{Tiny-jubjub} в её представлении Эдвардса (уравнение \ref{TJJ13-twisted-edwards}). 

Из \examplename{} \ref{TJJ13-twisted-edwards} мы знаем, что пара $(11,6)$ является правильной точкой на кривой \curvename{Tiny-jubjub}, и мы используем эту точку как вход для выполнения схемы. Мы получаем следующее:
\begin{center}
% https://www.graphviz.org/pdf/dotguide.pdf
\digraph[scale=0.4]{G2C}{
	forcelabels=true;
	//center=true;
	splines=ortho;
	nodesep= 2.0;
	n1 -> n4 [label="S_1=11" labeldistance="4" /*, color=lightgray */];
        n1 -> n4 /*[ color=lightgray ]*/
	n2 -> n5 [label="S_2=6" /*, color=lightgray */];
	n2 -> n5 /*[ color=lightgray ]*/ ;
	n3 -> n8 /*[ color=lightgray ]*/ ;
	n4 -> n8 [taillabel="S_3=4  " /*, color=lightgray */];
	n4 -> n10 [taillabel="S_3=4", labeldistance="4" /*, color=lightgray */];
	n5 -> n10 [xlabel="S_4=10 " /*, color=lightgray */];
	n5 -> n13 [headlabel="S_4=10", labeldistance="4" /*, color=lightgray */];
	n6 -> n13 /*[ color=lightgray ]*/ ;
	n7 -> n11 /*[ color=lightgray ]*/ ;
	n8 -> n11 [headlabel="[10*4=1]    "];
	n9 -> n12 /*[ color=lightgray ]*/ ;
	n10 -> n12 [taillabel="S_5=1  " /*, color=lightgray */];
	n11 -> n14 [headlabel="[1+1=2]    "];
	n12 -> n14 [label="  [8*1=8]"];	
	n13 -> n15 [headlabel="    [10*12=3]"];
	n14 -> n15 [taillabel="   [2+8=10]"];
	n15 -> n16 [label="  S_6=0", labeldistance="2" /*, color=lightgray */];
	n1 [shape=box, label="x" /*, color=lightgray */];
	n2 [shape=box, label="y" /*, color=lightgray */];
	n3 [shape=box, label="10" /*, color=lightgray */];
	n4 [label="*" /*, color=lightgray */];
	n5 [label="*" /*, color=lightgray */];
	n6 [shape=box, label="12" /*, color=lightgray */];
	n7 [shape=box, label="1" /*, color=lightgray */];
	n8 [label="*" /*, color=lightgray */];
	n9 [shape=box, label="8" /*, color=lightgray */];
	n10 [label="*" /*, color=lightgray */];
	n11 [label="+" /*, color=lightgray */];
	n12 [label="*" /*, color=lightgray */];	
	n13 [label="*" /*, color=lightgray */];
	n14 [label="+" /*, color=lightgray */];
	n15 [label="+" /*, color=lightgray */];
	n16 [shape=box, label="0" /*, color=lightgray */];		
}
\end{center}
Выполняя схему, мы действительно вычисляем $S_6=0$, как и ожидалось, что доказывает, что $(11,6)$ является точкой на кривой \curvename{Tiny-jubjub} в её представлении Эдвардса. Корректное присваивание $C_{tiny-jj}(\F_{13})$ поэтому даётся следующей строкой:
$$
S_{tiny-jj} = <S_1, S_2, S_3, S_4, S_5, S_6> = <11, 6, 4, 10, 1, 0>
$$
\end{example}
\subsubsection{Выполнимость схемы}
\label{circuit-satisfiability} Чтобы понять, как алгебраические схемы порождают формальные языки, заметим, что каждая алгебраическая схема $C(\F)$ над полем $\F$ определяет функцию решения над алфавитом $\Sigma_I \times \Sigma_W = \F \times \F$ следующим образом:
\begin{equation}
R_{C(\F)} : \F^* \times \F^* \to \{true, false\}\;;\;
(I;W) \mapsto
\begin{cases}
true & (I;W) \text{ является корректным присваиванием для } C(\F)\\
false & \text{ иначе}
\end{cases}
\end{equation}
Каждая алгебраическая схема поэтому определяет формальный язык. Грамматика этого языка закодирована в форме схемы, слова — это присваивания меткам рёбер, выведенные из выполнения схемы, а \term{утверждения} (statements) — это утверждения о знании ``Учитывая экземпляр $I$, существует свидетель $W$, такой что $(I;W)$ является корректным присваиванием схеме''. Конструктивное доказательство этого утверждения поэтому является присваиванием элемента поля каждой переменной свидетеля, что верифицируется выполнением схемы, чтобы увидеть, совпадает ли присваивание выполнения с присваиванием доказательства.

В контексте систем доказательств с нулевым разглашением выполнение схем также часто называется \term{генерацией свидетеля} (witness generation), поскольку в приложениях часть экземпляра обычно является публичной, в то время как задача доказывающего — вычислить часть свидетеля.

\begin{remark}[Выполнимость схемы] Подобно \ref{r1cs-satisfiability}, следует отметить, что в нашем определении каждая схема определяет свой собственный язык. Однако в более теоретических подходах часто рассматривается другой язык, обычно называемый \term{выполнимостью схемы} (circuit satisfiability), который полезен, когда речь идёт о более абстрактных проблемах, таких как выразительность или вычислительная сложность класса \term{всех} алгебраических схем над данным полем. С нашей точки зрения, язык выполнимости схемы получается объединением всех языков схем, которые есть в нашем определении. Будем более точны. Пусть алфавит $\Sigma=\F$ — поле. Тогда
$$
L_{CIRCUIT\_SAT(\F)} = \{(i;w)\in \Sigma^*\times \Sigma^*\;|\; \text{существует схема } C(\F) \text{ такая, что } (i;w) \text{ является корректным присваиванием}\}
$$
\end{remark}
\begin{example}[3-Factorization]Рассмотрим снова схему $C_{3.fac}$ из \examplename{} \ref{ex:3-fac-zk-circuit}. Мы называем ассоциированный язык $L_{3.fac\_circ}$.

Чтобы понять, как выглядит конструктивное доказательство утверждения в $L_{3.fac\_circ}$, рассмотрим экземпляр $I_1= 11$. Чтобы предоставить доказательство утверждения ``Существует свидетель $W$, такой что $(I_1;W)$ является словом в $L_{3.fac\_circ}$'' доказательство поэтому должно состоять из корректных значений переменных $W_1$, $W_2$, $W_3$ и $W_4$. Любой доказывающий поэтому должен найти входные значения для $W_1$, $W_2$ и $W_3$, а затем выполнить схему, чтобы вычислить $W_4$ при предположении $I_1=11$. 

Пример \ref{ex:3-fac-zk-circuit_2} подразумевает, что $<2,3,4,6>$ является правильным конструктивным доказательством. Чтобы верифицировать доказательство, верификатору нужно выполнить схему с экземпляром $I_1=11$ и входами $W_1=2$, $W_2=3$ и $W_3=4$, чтобы решить, является ли доказательство корректным присваиванием или нет. 
\end{example}
\begin{exercise}
Рассмотрим схему $C_{tiny-jj}(\F_{13})$ из \examplename{} \ref{ex:TJJ-circuit_1}, с ассоциированным языком $L_{tiny-jj}$. Сконструируйте доказательство $\pi$ для экземпляра $<11,6>$ и верифицируйте доказательство.
\end{exercise}

\subsubsection{Ассоциированные системы ограничений}
\label{sec:circuits_associated_R1CS} Как мы видели в \ref{sec:R1CS}, \concept{системы ограничений ранга 1} определяют способ представления утверждений в виде системы квадратичных уравнений над конечными полями, подходящих для основанных на спаривании систем доказательства с нулевым разглашением. Однако эти уравнения не предоставляют практического способа для доказывающего фактически вычислить решение. С другой стороны, алгебраические схемы могут быть выполнены, чтобы эффективно вывести корректные присваивания.

В этом параграфе мы покажем, как преобразовать любую алгебраическую схему в \concept{систему ограничений ранга 1} таким образом, что корректные присваивания схемы находятся во взаимно однозначном соответствии с решениями ассоциированной R1CS.

Чтобы увидеть это, пусть $C(\F)$ — алгебраическая схема над конечным полем $\F$, со строкой меток рёбер $<S_1,S_2,\ldots, S_n>$. Тогда мы начинаем с пустой R1CS и для каждой метки ребра $S_j$ из этого множества выполняется один из следующих шагов:
\begin{itemize}
\label{algebraic-gates}
\item Если метка ребра $S_j$ является исходящим ребром вентиля умножения, R1CS получает новое квадратичное ограничение
\begin{equation}
(\text{левый вход})\cdot (\text{правый вход}) = S_j
\end{equation} 
В этом выражении $(\text{левый вход})$ — это выход из символьного выполнения подграфа, который состоит из левого входного ребра этого вентиля и всех рёбер и узлов, которые имеют это ребро на своём пути, начиная с константных входов или помеченных исходящих рёбер других узлов.

Аналогичным образом $(\text{правый вход})$ — это выход из символьного выполнения подграфа, который состоит из правого входного ребра этого вентиля и всех рёбер и узлов, которые имеют это ребро на своём пути, начиная с константных входов или помеченных исходящих рёбер других узлов.
\item Если метка ребра $S_j$ является исходящим ребром вентиля сложения, R1CS получает новое квадратичное ограничение
\begin{equation}
(\text{левый вход} + \text{правый вход})\cdot 1 = S_j
\end{equation}  
В этом выражении $(\text{левый вход})$ — это выход из символьного выполнения подграфа, который состоит из левого входного ребра этого вентиля и всех рёбер и узлов, которые имеют это ребро на своём пути, начиная с константных входов или помеченных исходящих рёбер других узлов.

Аналогичным образом $(\text{правый вход})$ — это выход из символьного выполнения подграфа, который состоит из правого входного ребра этого вентиля и всех рёбер и узлов, которые имеют это ребро на своём пути, начиная с константных входов или помеченных исходящих рёбер других узлов.
\item Никакая другая метка ребра не добавляет ограничение в систему.
\end{itemize}

Если алгебраическая схема $C(\F)$ построена согласно правилам из \ref{def:algebraic-circuit}, результат этого метода — \concept{система ограничений ранга 1}, и, в этом смысле, каждая алгебраическая схема $C(\F)$ порождает R1CS $R$, которую мы называем \term{ассоциированной R1CS} схемы. Можно показать, что строка элементов поля $<S_1,S_2,\ldots, S_n>$ является корректным присваиванием схеме тогда и только тогда, когда та же строка является решением ассоциированной R1CS. Выполнения схем поэтому эффективно вычисляют решения \concept{систем ограничений ранга 1}.

Чтобы понять вклад арифметических вентилей в число ограничений, заметим, что,
согласно построению \ref{def:algebraic-circuit}, вентили умножения имеют метки на своих исходящих рёбрах тогда и только тогда, когда существует хотя бы одно помеченное ребро в обоих входных путях, или если исходящее ребро является входом в стоковый узел. Это подразумевает, что умножение на константу по существу бесплатно в том смысле, что оно не добавляет нового ограничения в систему, пока этот вентиль умножения не является входом в выходной узел.

Более того, вентили сложения имеют метки на своих исходящих рёбрах тогда и только тогда, когда они являются входами в стоковые узлы. Это подразумевает, что сложение по существу бесплатно в том смысле, что оно не добавляет нового ограничения в систему, пока этот вентиль сложения не является входом в выходной узел.

\begin{example}[$3$-factorization] Рассмотрим нашу проблему $3$-факторизации из \examplename{} \ref{ex:L-3fac-zk} и ассоциированную схему $C_{3.fac}(\F_{13})$ из \examplename{} \ref{ex:3-fac-zk-circuit}. Наша задача — преобразовать эту схему в эквивалентную \concept{систему ограничений ранга 1}.
\begin{center}
\digraph[scale=0.5]{G1}{
	forcelabels=true;
	//center=true;
	splines=ortho;
	nodesep= 2.0;
	n1 -> n3 [xlabel="W_2  "];
	n2 -> n3 [xlabel="  W_1"];
	n3 -> n5 [label="W_4"];
	n4 -> n5 [label=" W_3"];
	n5 -> n6 [label=" I_1"];
	n1 [shape=box, label="x_2"];
	n2 [shape=box, label="x_1"];
	n3 [label="*"];
	n4 [shape=box, label="x_3"];
	n5 [label="*"];
	n6 [shape=box, label="f_(3.fac_zk)"];
}
\end{center}
Мы начинаем с пустой R1CS, и, чтобы сгенерировать все ограничения, мы должны итерировать по множеству меток рёбер $<I_1;W_1,W_2,W_3,W_4>$.  

Начиная с метки ребра $I_1$, мы видим, что это исходящее ребро вентиля умножения, и, поскольку оба входных ребра помечены, мы должны добавить следующее ограничение в систему:
\begin{align*}
(\text{левый вход})\cdot (\text{правый вход})  &= I_1 & \Leftrightarrow\\
W_4\cdot W_3  &= I_1
\end{align*}
Далее мы рассматриваем метку ребра $W_1$ и, поскольку она не является исходящим ребром вентиля умножения или сложения, мы не добавляем ограничение в систему. То же самое верно для меток $W_2$ и $W_3$.

Для метки ребра $W_4$ мы видим, что это исходящее ребро вентиля умножения, и, поскольку оба входных ребра помечены, мы должны добавить следующее ограничение в систему:
\begin{align*}
(\text{левый вход})\cdot (\text{правый вход})  &= W_4 & \Leftrightarrow\\
W_2\cdot W_1  &= W_4
\end{align*} 
Поскольку больше нет помеченных рёбер, все ограничения сгенерированы, и мы должны объединить их, чтобы получить ассоциированную R1CS $C_{3.fac}(\F_{13})$:
\begin{align*}
 W_4\cdot W_3 & = I_1\\
 W_2\cdot W_1 & = W_4
\end{align*}
Эта система эквивалентна R1CS, которую мы вывели в \examplename{} \ref{ex:3-factorization-r1cs}. Языки $L_{3.fac\_zk}$ и $L_{3.fac\_circ}$ поэтому эквивалентны, и как схема, так и R1CS — просто два разных способа выражения одного и того же языка.
\end{example}
\begin{example}
\label{ex:tjj-circuit-2-tjj-R1CS} Чтобы рассмотреть более общее преобразование, мы снова рассмотрим схему \curvename{Tiny-jubjub} из \examplename{} \ref{ex:TJJ-circuit_2}. Правильная схема даётся следующим графом, где мы выделили все узлы, которые вносят ограничение в R1CS:
\begin{center}
\digraph[scale=0.4]{G2D}{
	forcelabels=true;
	center=true;
	splines=ortho;
	nodesep= 2.0;
	n1 -> n4 [label="S_1" labeldistance="4" /*, color=lightgray */];
        n1 -> n4 /*[ color=lightgray ]*/
	n2 -> n6 [label="S_2" /*, color=lightgray */];
	n2 -> n6 /*[ color=lightgray ]*/ ;
	n3 -> n9 /*[ color=lightgray ]*/ ;
	n4 -> n9 [taillabel="S_3", labeldistance="2" /*, color=lightgray */];
	n4 -> n13 [taillabel="S_3", labeldistance="4" /*, color=lightgray */];
	n5 -> n10 /*[ color=lightgray ]*/ ;
	n6 -> n10 [headlabel="S_4", labeldistance="4" /*, color=lightgray */];
	n6 -> n11 [taillabel="S_4", labeldistance="4" /*, color=lightgray */];
	n7 -> n13 /*[ color=lightgray ]*/ ;
	n8 -> n12 /*[ color=lightgray ]*/ ;
	n9 -> n12 /*[ color=lightgray ]*/ ;
	n10 -> n13 /*[ color=lightgray ]*/ ; 
	n11 -> n15 /*[ color=lightgray ]*/ ;
	n12 -> n14 /*[ color=lightgray ]*/ ;	
	n13 -> n14 [xlabel="S_5  "  /*, color=lightgray */];
	n14 -> n15 /*[ color=lightgray ]*/ ;
	n15 -> n16 [label="  S_6=0", labeldistance="2" /*, color=lightgray */];
	n1 [shape=box, label="x" /*, color=lightgray */];
	n2 [shape=box, label="y" /*, color=lightgray */];
	n3 [shape=box, label="10" /*, color=lightgray */];
	n4 [label="*", style=filled /*, color=lightgray */];
	n5 [shape=box, label="8" /*, color=lightgray */];
	n6 [label="*", style=filled /*, color=lightgray */];
	n7 [shape=box, label="12" /*, color=lightgray */];
	n8 [shape=box, label="1" /*, color=lightgray */];
	n9 [label="*" /*, color=lightgray */];
	n10 [label="*" /*, color=lightgray */];
	n11 [label="*" /*, color=lightgray */];	
	n12 [label="+" /*, color=lightgray */];	
	n13 [label="*", style=filled /*, color=lightgray */];
	n14 [label="+" /*, color=lightgray */];
	n15 [label="+", style=filled /*, color=lightgray */];
	n16 [shape=box, label="0" /*, color=lightgray */];		
}
\end{center}
Чтобы вычислить число ограничений, заметим, что у нас есть $3$ вентиля умножения, которые имеют метки на своих исходящих рёбрах, и $1$ вентиль сложения, который имеет метку на своём исходящем ребре. Мы поэтому должны вычислить $4$ квадратичных ограничения.

Чтобы вывести ассоциированную R1CS, мы должны начать с пустой R1CS, а затем итерировать по множеству $<S_1,S_2,S_3,S_4,S_5,S_6=0>$ всех меток рёбер, чтобы сгенерировать ограничения.

Рассматривая метку ребра $S_1$, мы видим, что ассоциированные рёбра не являются исходящими рёбрами какого-либо арифметического вентиля, и поэтому мы не должны добавлять новое ограничение в систему. То же самое верно для метки ребра $S_2$. Глядя на метку ребра $S_3$, мы видим, что ассоциированные рёбра являются исходящими рёбрами вентиля умножения, и что ассоциированный подграф даётся:
\begin{center}
\digraph[scale=0.4]{G2E}{
	forcelabels=true;
	center=true;
	splines=ortho;
	nodesep= 2.0;
	n1 -> n4 [label="S_1" labeldistance="4" /*, color=lightgray */];
    n1 -> n4 /*[ color=lightgray ]*/
	n2 -> n6 [label="S_2", color=lightgray];
	n2 -> n6 [ color=lightgray ] ;
	n3 -> n9 [ color=lightgray ] ;
	n4 -> n9 [taillabel="S_3", labeldistance="2" /*, color=lightgray */];
	n4 -> n13 [taillabel="S_3", labeldistance="4" /*, color=lightgray */];
	n5 -> n10 [ color=lightgray ] ;
	n6 -> n10 [headlabel="S_4", labeldistance="4", color=lightgray];
	n6 -> n11 [taillabel="S_4", labeldistance="4", color=lightgray];
	n7 -> n13 [ color=lightgray ] ;
	n8 -> n12 [ color=lightgray ] ;
	n9 -> n12 [ color=lightgray ] ;
	n10 -> n13 [ color=lightgray ] ; 
	n11 -> n15 [ color=lightgray ] ;
	n12 -> n14 [ color=lightgray ] ;	
	n13 -> n14 [xlabel="S_5  ", color=lightgray];
	n14 -> n15 [ color=lightgray ] ;
	n15 -> n16 [label="  S_6=0", labeldistance="2", color=lightgray];
	n1 [shape=box, label="x" /*, color=lightgray */];
	n2 [shape=box, label="y", color=lightgray];
	n3 [shape=box, label="10", color=lightgray ];
	n4 [label="*", style=filled /*, color=lightgray */];
	n5 [shape=box, label="8", color=lightgray];
	n6 [label="*", style=filled, color=lightgray];
	n7 [shape=box, label="12", color=lightgray];
	n8 [shape=box, label="1", color=lightgray];
	n9 [label="*", color=lightgray];
	n10 [label="*", color=lightgray];
	n11 [label="*", color=lightgray];	
	n12 [label="+", color=lightgray];	
	n13 [label="*", style=filled, color=lightgray];
	n14 [label="+", color=lightgray];
	n15 [label="+", style=filled, color=lightgray];
	n16 [shape=box, label="0", color=lightgray];		
}
\end{center}
И левый, и правый вход этого вентиля умножения помечены $S_1$. Мы поэтому должны добавить следующее ограничение в систему:
$$
S_1 \cdot S_1 = S_3
$$
Глядя на метку ребра $S_4$, мы видим, что ассоциированные рёбра являются исходящими рёбрами вентиля умножения, и что ассоциированный подграф даётся:
\begin{center}
\digraph[scale=0.4]{G2F}{
	forcelabels=true;
	center=true;
	splines=ortho;
	nodesep= 2.0;
	n1 -> n4 [label="S_1" labeldistance="4", color=lightgray];
        n1 -> n4 [ color=lightgray ]
	n2 -> n6 [label="S_2" /*, color=lightgray */];
	n2 -> n6 /* [ color=lightgray ] */ ;
	n3 -> n9 [ color=lightgray ] ;
	n4 -> n9 [taillabel="S_3", labeldistance="2", color=lightgray];
	n4 -> n13 [taillabel="S_3", labeldistance="4", color=lightgray];
	n5 -> n10 [ color=lightgray ] ;
	n6 -> n10 [headlabel="S_4", labeldistance="4" /*, color=lightgray */];
	n6 -> n11 [taillabel="S_4", labeldistance="4" /*, color=lightgray */];
	n7 -> n11 [ color=lightgray ] ;
	n8 -> n12 [ color=lightgray ] ;
	n9 -> n12 [ color=lightgray ] ;
	n10 -> n13 [ color=lightgray ] ; 
	n11 -> n15 [ color=lightgray ] ;
	n12 -> n14 [ color=lightgray ] ;	
	n13 -> n14 [xlabel="S_5  ", color=lightgray];
	n14 -> n15 [ color=lightgray ] ;
	n15 -> n16 [label="  S_6=0", labeldistance="2", color=lightgray];
	n1 [shape=box, label="x", color=lightgray];
	n2 [shape=box, label="y", /* color=lightgray */];
	n3 [shape=box, label="10", color=lightgray];
	n4 [label="*", style=filled, color=lightgray];
	n5 [shape=box, label="8", color=lightgray];
	n6 [label="*", style=filled /*, color=lightgray */];
	n7 [shape=box, label="12", color=lightgray];
	n8 [shape=box, label="1", color=lightgray];
	n9 [label="*", color=lightgray];
	n10 [label="*", color=lightgray];
	n11 [label="*", color=lightgray];	
	n12 [label="+", color=lightgray];	
	n13 [label="*", style=filled, color=lightgray];
	n14 [label="+", color=lightgray];
	n15 [label="+", style=filled, color=lightgray];
	n16 [shape=box, label="0", color=lightgray];		
}
\end{center}
И левый, и правый вход этого вентиля умножения помечены $S_2$, и мы поэтому должны добавить следующее ограничение в систему:
$$
S_2 \cdot S_2 = S_4
$$
Метка ребра $S_5$ более интересна. Чтобы увидеть, как она подразумевает ограничение, мы должны сначала сконструировать ассоциированный подграф, который состоит из всех рёбер, узлов и путей, начинающихся либо с константного входа, либо с помеченного ребра. Мы получаем
\begin{center}
\digraph[scale=0.4]{G2G}{
	forcelabels=true;
	center=true;
	splines=ortho;
	nodesep= 2.0;
	n1 -> n4 [label="S_1" labeldistance="4", color=lightgray];
        n1 -> n4 [ color=lightgray ]
	n2 -> n6 [label="S_2", color=lightgray];
	n2 -> n6 [ color=lightgray ] ;
	n3 -> n9 [ color=lightgray ] ;
	n4 -> n9 [taillabel="S_3", labeldistance="2", color=lightgray];
	n4 -> n13 [taillabel="S_3", labeldistance="4" /*, color=lightgray */];
	n5 -> n10 /*[ color=lightgray ]*/ ;
	n6 -> n10 [headlabel="S_4", labeldistance="4" /*, color=lightgray */];
	n6 -> n11 [taillabel="S_4", labeldistance="4", color=lightgray ];
	n7 -> n11 [ color=lightgray ] ;
	n8 -> n12 [ color=lightgray ] ;
	n9 -> n12 [ color=lightgray ] ;
	n10 -> n13 /*[ color=lightgray ]*/ ; 
	n11 -> n15 [ color=lightgray ] ;
	n12 -> n14 [ color=lightgray ] ;	
	n13 -> n14 [xlabel="S_5  "  /*, color=lightgray */];
	n14 -> n15 [ color=lightgray ] ;
	n15 -> n16 [label="  S_6=0", labeldistance="2" , color=lightgray ];
	n1 [shape=box, label="x", color=lightgray];
	n2 [shape=box, label="y", color=lightgray];
	n3 [shape=box, label="10", color=lightgray];
	n4 [label="*", style=filled /*, color=lightgray*/];
	n5 [shape=box, label="8" /*, color=lightgray */];
	n6 [label="*", style=filled /*, color=lightgray */];
	n7 [shape=box, label="12", color=lightgray];
	n8 [shape=box, label="1", color=lightgray];
	n9 [label="*", color=lightgray];
	n10 [label="*" /*, color=lightgray */];
	n11 [label="*", color=lightgray];	
	n12 [label="+", color=lightgray];	
	n13 [label="*", style=filled /*, color=lightgray */];
	n14 [label="+", color=lightgray];
	n15 [label="+", style=filled, color=lightgray];
	n16 [shape=box, label="0", color=lightgray];		
}
\end{center}
Правый вход ассоциированного вентиля умножения даётся помеченным ребром $S_3$. Однако левый вход не является помеченным ребром, но имеет помеченное ребро в одном из своих путей. Чтобы вычислить левый множитель этого ограничения, мы должны вычислить выход подграфа, ассоциированного с левым ребром, который равен $S_4\cdot 8$. Это даёт ограничение
$$
(S_4 \cdot 8) \cdot S_3 = S_5
$$ 
Последняя метка ребра — константа $S_6=0$. Чтобы увидеть, как она подразумевает ограничение, мы должны сконструировать ассоциированный подграф, который состоит из всех рёбер, узлов и путей, начинающихся либо с константного входа, либо с помеченного ребра. Мы получаем
\begin{center}
\digraph[scale=0.4]{G2H}{
	forcelabels=true;
	center=true;
	splines=ortho;
	nodesep= 2.0;
	n1 -> n4 [label="S_1" /* comment*/ labeldistance="4", color=lightgray];
    n1 -> n4 [ color=lightgray ]
	n2 -> n6 [label="S_2", color=lightgray];
	n2 -> n6 [ color=lightgray ] ;
	n3 -> n9 /*[ color=lightgray ]*/ ;
	n4 -> n9 [taillabel="S_3", labeldistance="2" /*, color=lightgray */];
	n4 -> n13 [taillabel="S_3", labeldistance="4" , color=lightgray ];
	n5 -> n10 [ color=lightgray ] ;
	n6 -> n10 [headlabel="S_4", labeldistance="4" , color=lightgray ];
	n6 -> n11 [taillabel="S_4", labeldistance="4" /*, color=lightgray */];
	n7 -> n11 /*[ color=lightgray ]*/ ;
	n8 -> n12 /*[ color=lightgray ]*/ ;
	n9 -> n12 /*[ color=lightgray ]*/ ;
	n10 -> n13 [ color=lightgray ] ; 
	n11 -> n15 /*[ color=lightgray ]*/ ;
	n12 -> n14 /*[ color=lightgray ]*/ ;	
	n13 -> n14 [xlabel="S_5  "  /*, color=lightgray */];
	n14 -> n15 /*[ color=lightgray ]*/ ;
	n15 -> n16 [label="  S_6=0", labeldistance="2" /*, color=lightgray */];
	n1 [shape=box, label="x", color=lightgray];
	n2 [shape=box, label="y", color=lightgray];
	n3 [shape=box, label="10" /*, color=lightgray */];
	n4 [label="*", style=filled /*, color=lightgray */];
	n5 [shape=box, label="8", color=lightgray];
	n6 [label="*", style=filled /*, color=lightgray */];
	n7 [shape=box, label="12" /*, color=lightgray */];
	n8 [shape=box, label="1" /*, color=lightgray */];
	n9 [label="*" /*, color=lightgray */];
	n10 [label="*" , color=lightgray];
	n11 [label="*" /*, color=lightgray */];	
	n12 [label="+" /*, color=lightgray */];	
	n13 [label="*", style=filled /*, color=lightgray */];
	n14 [label="+" /*, color=lightgray */];
	n15 [label="+", style=filled /*, color=lightgray */];
	n16 [shape=box, label="0" /*, color=lightgray */];		
}
\end{center}
И левый, и правый вход не помечены, но имеют помеченные рёбра в своём пути. Поскольку вентиль является вентилем сложения, правый множитель в квадратичном ограничении всегда $1$, а левый множитель вычисляется символьным выполнением всех входов во все вентили в подсхеме. Мы получаем
$$
(12\cdot S_4 + S_5 + 10\cdot S_3 + 1)\cdot 1 = 0
$$

Поскольку больше нет помеченных исходящих рёбер, мы закончили вывод ограничений. Объединяя все ограничения вместе, мы получаем следующую R1CS:
\begin{align*}
 S_1 \cdot S_1 &= S_3\\
 S_2 \cdot S_2 &= S_4\\
 (S_4\cdot 8)\cdot S_3 &= S_5\\
 (12\cdot S_4 + S_5 + 10\cdot S_3 + 1)\cdot 1 &= 0
\end{align*}
что эквивалентно R1CS, которую мы вывели в \examplename{} \ref{ex:TJJ-r1cs}. Таким образом, как схема, так и R1CS — просто два разных способа выразить один и тот же язык.
\end{example}
\subsection{Квадратичные арифметические программы}
\label{sec:QAP}
 Мы ввели алгебраические схемы и их ассоциированные \concept{системы ограничений ранга 1} как две особые модели, способные представлять ограниченные вычисления. Обе модели определяют формальные языки, и ассоциированные утверждения о принадлежности, а также утверждения о знании могут быть доказаны конструктивным образом путём выполнения схемы, чтобы вычислить решения её ассоциированной R1CS.

Одна из причин, по которым эти системы полезны в контексте кратких систем доказательства с нулевым разглашением, заключается в том, что любая R1CS может быть преобразована в другую вычислительную модель, называемую \term{квадратичной арифметической программой} (Quadratic Arithmetic Program, QAP), которая служит основой для некоторых из наиболее эффективных кратких неинтерактивных генераторов доказательств с нулевым разглашением, которые в настоящее время существуют.

Как мы увидим, доказательство утверждений для языков, у которых функции решения определены квадратичными арифметическими программами, может быть достигнуто предоставлением определённых многочленов, и эти доказательства могут быть верифицированы проверкой определённого свойства делимости этих многочленов.
 
\subsubsection{Представление QAP} Чтобы понять, что такое квадратичные арифметические программы подробно, пусть $\F$ — поле и $R$ — \concept{система ограничений ранга 1} над $\F$, такая что число ненулевых элементов в $\F$ строго больше числа $k$ ограничений в $R$. Более того, пусть $a_j^i$, $b_j^i$ и $c_j^i\in\F$ для каждого индекса $0\leq j \leq n+m$ и $1\leq i \leq k$, являются определяющими константами R1CS, и $m_1$, $\ldots$, $m_k$ — произвольные, обратимые и различные элементы из $\F$.
  
Тогда \term{квадратичная арифметическая программа}, ассоциированная с R1CS $R$, — это следующее множество многочленов над $\F$:
\begin{equation}
\label{def:QAP-target-poly}
QAP(R) = \left\{T\in \F[x],\left\{A_j,B_j,C_j\in \F[x]\right\}_{j=0}^{n+m}\right\}
\end{equation}
Здесь $T(x) := \Pi_{l=1}^k (x- m_l)$ — многочлен степени $k$, называемый \term{целевым многочленом} (target polynomial) QAP, а $A_j$, $B_j$, а также $C_j$ — единственные многочлены степени $k-1$, определяемые следующим уравнением:
\begin{equation}
\label{def:QAP-polynomials}
\begin{array}{lllr}
A_j(m_i)=a_j^i, & B_j(m_i)=b_j^i, & C_j(m_i)=c_j^i & \text{ для всех } j= 1, \ldots , n+m+1, i=1,\ldots,k 
\end{array}
\end{equation}
Учитывая некоторую \concept{систему ограничений ранга 1}, ассоциированная квадратичная арифметическая программа поэтому является множеством многочленов, вычисленных из констант в R1CS, вместе с целевым многочленом. Заметим, что целевой многочлен обращается в $0$ для каждого $m_1$, $\ldots$, $m_k$, что является полезным фактом для верификации доказательств. Чтобы увидеть, что многочлены $A_j$, $B_j$ и $C_j$ однозначно определяются уравнениями \ref{def:QAP-polynomials}, вспомним, что многочлен степени $k-1$ полностью определяется $k$ точками оценки, и он может быть вычислен, например, интерполяцией Лагранжа, описанной выше в Алгоритме \ref{alg_lagrange_interplation}.

Вычисление QAP из любой данной R1CS может быть достигнуто в следующих трёх шагах. Если R1CS состоит из $k$ ограничений, сначала выберем $k$ различных, обратимых элементов из поля $\F$. Каждый выбор определяет различную QAP для той же R1CS. Затем вычислим целевой многочлен $T$ согласно его определению в \ref{def:QAP-target-poly}. После этого используем интерполяцию Лагранжа через Алгоритм \ref{alg_lagrange_interplation}, чтобы вычислить многочлены $A_j$ для каждого $1\leq j \leq k$ из множества
\begin{equation}
S_{A_j} = \{(m_1,a^1_j),\ldots,(m_k,a^k_j)\}
\end{equation}
Затем используем тот же метод, чтобы вычислить многочлены $B_j$ и $C_j$ для каждого $1\leq j \leq k$.
\begin{example}[3-factorization]
\label{ex:3-fac-QAP}  Чтобы дать лучшую интуицию о квадратичных арифметических программах и о том, как они вычисляются из их ассоциированных \concept{систем ограничений ранга 1}, рассмотрим язык $L_{3.fac\_zk}$ из \examplename{} \ref{ex:L-3fac-zk} и его ассоциированную R1CS из \examplename{} \ref{ex:3-factorization-r1cs}:
\begin{align*}
W_1 \cdot W_2 & = W_4 & \text{ограничение } 1\\
W_4 \cdot W_3 & = I_1 & \text{ограничение } 2
\end{align*}
В этом примере мы хотим преобразовать эту R1CS в ассоциированную QAP. Согласно \examplename{} \ref{ex:3-factorization-r1cs} определяющие константы $a_j^i$, $b_j^i$ и $c_j^i$ R1CS даются следующим образом:
$$
\begin{array}{llllll}
a_0^1 = 0 & a_1^1= 0 & a_2^1= 1 & a_3^1 = 0 & a_4^1= 0  & a_5^1= 0 \\ 
a_0^2 = 0 & a_1^2= 0 & a_2^2= 0 & a_3^2 = 0 & a_4^2= 0  & a_5^2= 1 \\ 
\\
b_0^1 = 0 & b_1^1= 0 & b_2^1= 0 & b_3^1 = 1 & b_4^1= 0  & b_5^1= 0 \\ 
b_0^2 = 0 & b_1^2= 0 & b_2^2= 0 & b_3^2 = 0 & b_4^2= 1  & b_5^2= 0 \\ 
\\
c_0^1 = 0 & c_1^1= 0 & c_2^1= 0 & c_3^1 = 0 & c_4^1= 0  & c_5^1= 1 \\ 
c_0^2 = 0 & c_1^2= 1 & c_2^2= 0 & c_3^2 = 0 & c_4^2= 0  & c_5^2= 0 
\end{array} 
$$
Поскольку R1CS определена над полем $\F_{13}$ и имеет два ограничения, нам нужно выбрать два произвольных, но обратимых и различных элемента $m_1$ и $m_2$ из $\F_{13}$. Мы выбираем $m_{1}=5$ и $m_{2}=7$, и с этим выбором мы получаем целевой многочлен
\begin{align*}
T(x) & = (x-m_1)(x-m_2) & \text{\# Definition of T}\\
     & = (x-5)(x-7)  & \text{\# Insert our chosen values}\\
     & = (x+8)(x+6)  & \text{\# Additive inverses in } \F_{13}\\
     & = x^2 + x +9 & \text{\# Expand}
\end{align*}
Затем мы должны вычислить многочлены $A_j$, $B_j$ и $C_j$ по их определяющему уравнению из коэффициентов R1CS. Поскольку R1CS имеет два ограничивающих уравнения, эти многочлены степени $1$ и определяются их значениями в точках $m_1=5$ и $m_2=7$.

В точке $m_1$ каждый многочлен $A_j$ определён как $a_j^1$, а в точке $m_2$ каждый многочлен $A_j$ определён как $a_j^2$. То же самое верно для многочленов $B_j$, а также $C_j$. Записывая все эти уравнения, мы получаем:
$$
\begin{array}{llllll}
A_0(5)=0, & A_1(5)=0, & A_2(5)=1, & A_3(5)=0, & A_4(5)=0, & A_5(5)=0 \\
A_0(7)=0, & A_1(7)=0, & A_2(7)=0, & A_3(7)=0, & A_4(7)=0, & A_5(7)=1\\
\\
B_0(5)=0, & B_1(5)=0, & B_2(5)=0, & B_3(5)=1, & B_4(5)=0, & B_5(5)=0 \\
B_0(7)=0, & B_1(7)=0, & B_2(7)=0, & B_3(7)=0, & B_4(7)=1, & B_5(7)=0\\
\\
C_0(5)=0, & C_1(5)=0, & C_2(5)=0, & C_3(5)=0, & C_4(5)=0, & C_5(5)=1 \\
C_0(7)=0, & C_1(7)=1, & C_2(7)=0, & C_3(7)=0, & C_4(7)=0, & C_5(7)=0
\end{array}
$$
Основная теорема алгебры подразумевает, что многочлен степени $k$, который равен нулю на $k+1$ точках, должен быть нулевым многочленом. Поскольку наши многочлены степени $1$ и определены на $2$ точках, мы поэтому знаем, что единственными ненулевыми многочленами в нашей QAP являются $A_2$, $A_5$, $B_3$, $B_4$, $C_1$ и $C_5$, и что мы можем использовать интерполяцию Лагранжа, чтобы вычислить их.

Чтобы вычислить $A_2$, мы замечаем, что множество $S_{A_2}$ в нашем примере даётся $S_{A_2}=\{(m_1,a^1_2), (m_2,a_2^2)\} = \{(5,1), (7,0)\}$. Используя это множество, мы получаем:
\begin{align*}
A_2(x) & = a^1_2\cdot \left(\frac{x-m_2}{m_1-m_2}\right) + a^2_2\cdot\left(\frac{x-m_1}{m_2-m_1}\right)
      = 1\cdot\left(\frac{x-7}{5-7}\right) + 0\cdot\left(\frac{x-5}{7-5}\right) \\
    & = \frac{x-7}{-2}
      = \frac{x-7}{11} & \text{\# } 11^{-1}=6 \\
    & = 6(x-7) 
      = 6x + 10 & \text{\# } -7 = 6 \text{ и } 6\cdot 6 = 10
\end{align*}
Чтобы вычислить $A_5$, мы замечаем, что множество $S_{A_5}$ в нашем примере даётся $S_{A_5}=\{(m_1,a^1_5), (m_2,a^2_5)\} = \{(5,0), (7,1)\}$. Используя это множество, мы получаем:
\begin{align*}
A_5(x) & = a^1_5\cdot\left(\frac{x-m_2}{m_1-m_2}\right) + a^2_5\cdot\left(\frac{x-m_1}{m_2-m_1}\right)
      = 0\cdot\left(\frac{x-7}{5-7}\right) + 1\cdot\left(\frac{x-5}{7-5}\right) \\
    & = \frac{x-5}{2} & \text{\# } 2^{-1}=7 \\
    & = 7(x-5) 
      = 7x + 4 & \text{\# } -5 = 8 \text{ и } 7\cdot 8 = 4
\end{align*}
Используя интерполяцию Лагранжа, мы можем заключить, что $A_2=B_3=C_5$, а также $A_5=B_4=C_1$, поскольку они являются многочленами степени $1$, которые принимают одинаковые значения в $2$ точках. Используя это, мы получаем следующее множество многочленов:
\begin{center}
\begin{tabular}{|l|l|l|}\hline 
$A_{0}(x)=0 $ &$ B_{0}(x)=0   $ & $C_{0}(x)=0$ \tabularnewline\hline 
$A_1(x)=0 $ &$ B_1(x)=0   $ & $C_1(x)=7x+4$ \tabularnewline\hline 
$A_2(x)=6x+10$ &$ B_2(x)=0$ & $C_2(x)=0$ \tabularnewline\hline 
$A_3(x)=0    $ &$ B_3(x)=6x+10$ & $C_3(x)=0$ \tabularnewline\hline 
$A_4(x)=0$ &$ B_4(x)=7x+4  $ & $C_4(x)=0$ \tabularnewline\hline 
$A_5(x)=7x+4$ &$ B_5(x)=0      $ & $C_5(x)=6x+10$ \tabularnewline\hline 
\end{tabular}
\end{center}
Мы можем использовать Sage, чтобы верифицировать наше вычисление. В Sage каждое кольцо многочленов имеет функцию \code{lagrange\_polynomial}, которая принимает определяющие точки как входы и возвращает ассоциированный многочлен Лагранжа.
\begin{sagecommandline}
sage: F13 = GF(13)
sage: F13t.<t> = F13[]
sage: T = F13t((t-5)*(t-7))
sage: A2 = F13t.lagrange_polynomial([(5,1),(7,0)])
sage: A5 = F13t.lagrange_polynomial([(5,0),(7,1)])
sage: T == F13t(t^2 + t + 9)
sage: A2 == F13t(6*t + 10)
sage: A5 == F13t(7*t + 4)
\end{sagecommandline}

Объединяя это вычисление с целевым многочленом, который мы вывели ранее, квадратичная арифметическая программа, ассоциированная с \concept{системой ограничений ранга 1} $R_{3.fac\_zk}$, даётся следующим образом:
\begin{multline}
\label{QAP-R3-fac-zk}
QAP(R_{3.fac\_zk}) =\{x^{2}+x+9,\notag\\
 \{0,0,6x+10,0,0,7x+4\},\{0,0,0,6x+10,7x+4,0\},\{0,7x+4,0,0,0,6x+10\}\}
\end{multline}
\end{example}
\begin{exercise}
Рассмотрим \concept{систему ограничений ранга 1} для точек на кривой \curvename{Tiny-jubjub} из \examplename{} \ref{ex:TJJ-r1cs}. Вычислите ассоциированную QAP для этой R1CS и перепроверьте своё вычисление с помощью sage.
\end{exercise}
\subsubsection{Выполнимость QAP} Одно из основных преимуществ квадратичных арифметических программ в системах доказательства заключается в том, что решения их ассоциированных \concept{систем ограничений ранга 1} находятся во взаимно однозначном соответствии с определёнными многочленами $P$, делящимися на целевой многочлен $T$ QAP. Верификация решений R1CS, а следовательно, проверка правильного выполнения схемы затем достижима делением многочлена $P$ на $T$.

Будем более конкретны. Пусть $R$ — некоторая \concept{система ограничений ранга 1} с ассоциированными переменными $(<I_1,\ldots,I_n>; <W_1,\ldots, W_m>)$, и пусть $QAP(R)$ — квадратичная арифметическая программа $R$. Тогда строка $(<I_1,\ldots,I_n>; <W_1,\ldots, W_m>)$ является решением R1CS тогда и только тогда, когда следующий многочлен делится на целевой многочлен $T$:
\begin{equation}\label{polynomial-P-IW}
P_{(I;W)} = \scriptstyle \left(A_0 + \sum_{j}^n I_j\cdot A_j + \sum_{j}^m W_j\cdot A_{n+j} \right) \cdot \left(B_0 + \sum_{j}^n I_j\cdot B_j + \sum_{j}^m W_j\cdot B_{n+j} \right) 
-\left(C_0 + \sum_{j}^n I_j\cdot C_j + \sum_{j}^m W_j\cdot C_{n+j} \right)
\end{equation}
Это работает, потому что $P_{(I;W)}$ будет равен $0$ для всех значений $m_1$, $\ldots$, $m_k$ тогда и только тогда, когда строка $(<I_1,\ldots,I_n>; <W_1,\ldots, W_m>)$ является решением R1CS. Поскольку целевой многочлен $T$ равен $0$ для всех значений $m_1$, $\ldots$, $m_k$, $P_{(I;W)}$ должен быть кратен $T$, чтобы также равняться $0$. Таким образом, наше предложенное решение будет корректным тогда и только тогда, когда $P_{(I;W)}$ делится на $T$.

Чтобы понять, как квадратичные арифметические программы определяют формальные языки, заметим, что каждая QAP над полем $\F$ определяет функцию решения над алфавитом $\Sigma_I \times \Sigma_W = \F \times \F$ следующим образом:
\begin{equation}
R_{QAP} : (\F)^* \times (\F)^* \to \{true, false\}\;;\;
(I;W) \mapsto
\begin{cases}
true & P_{(I;W)} \text{ делится на } T\\
false & \text{ иначе}
\end{cases}
\end{equation}
Это означает, что каждая QAP определяет формальный язык $L_{QAP}$, и, если QAP ассоциирована с R1CS, язык, порождённый QAP, и язык, порождённый R1CS, эквивалентны. В контексте языков, порождённых квадратичными арифметическими программами, \term{утверждение} (statement) — это утверждение о принадлежности ``Существует слово $(I;W)$ в $L_{QAP}$''. Доказательство этого утверждения поэтому даётся многочленом $P_{(I;W)}$, который верифицируется делением $P_{(I;W)}$ на $T$.

Заметьте структурные сходства и различия в определении R1CS и её ассоциированного языка в \ref{R1CS}, схем и их ассоциированных языков в \ref{sec:circuits} и QAP и их ассоциированных языков, как объяснено в этой части. Для схем и их ассоциированных \concept{систем ограничений ранга 1} конструктивное доказательство состоит из корректного присваивания элементов поля рёбрам схемы, или переменным в R1CS. Однако в случае QAP корректное доказательство состоит из многочлена $P_{(I;W)}$.

Чтобы вычислить конструктивное доказательство для утверждения в $L_{QAP}$, учитывая некоторый экземпляр $I$, доказывающий сначала должен вычислить конструктивное доказательство $W$ ассоциированной R1CS, например, выполнив схему R1CS. Имея $(I;W)$, доказывающий может затем вычислить многочлен $P_{(I;W)}$ и опубликовать многочлен как доказательство.

Верификация конструктивного доказательства в случае схемы достигается выполнением схемы, а затем сравнением результата с данным доказательством. Верификация того же доказательства в картине R1CS означает проверку, удовлетворяют ли элементы доказательства уравнениям R1CS. Напротив, верификация доказательства в картине QAP выполняется делением многочлена доказательства $P$ на целевой многочлен $T$. Доказательство верифицировано тогда и только тогда, когда $P$ делится на $T$.

\begin{example} Рассмотрим квадратичную арифметическую программу $QAP(R_{3.fac\_zk})$ из \examplename{} \ref{ex:3-fac-QAP} и её ассоциированную R1CS из уравнения \ref{ex:3-factorization-r1cs}. Чтобы дать интуицию о том, как выглядят доказательства в языке $L_{QAP(R_{3.fac\_zk})}$, рассмотрим экземпляр $I_1=11$. Как мы знаем из \examplename{} \ref{ex:3-fac-zk-circuit_2}, $(W_1,W_2,W_3,W_4)=(2,3,4,6)$ является правильным свидетелем, поскольку
$(<I_1>;<W_1,W_2,W_3,W_4>)=(<11>;<2,3,4,6>)$ является корректным присваиванием схемы, а следовательно, решением $R_{3.fac\_zk}$ и конструктивным доказательством для языка $L_{R_{3.fac\_zk}}$.

Чтобы преобразовать это конструктивное доказательство в доказательство знания в языке $L_{QAP(R_{3.fac\_zk})}$, доказывающий должен использовать элементы конструктивного доказательства, чтобы вычислить многочлен $P_{(I;W)}$.

В случае $(<I_1>;<W_1,W_2,W_3,W_4>)=(<11>;<2,3,4,6>)$ ассоциированное доказательство вычисляется следующим образом:
\begin{align*}
P_{(I;W)}  = & \scriptstyle \left(A_0 + \sum_{j}^n I_j\cdot A_j + \sum_{j}^m W_j\cdot A_{n+j} \right) \cdot \left(B_0 + \sum_{j}^n I_j\cdot B_j + \sum_{j}^m W_j\cdot B_{n+j} \right) 
-\left(C_0 + \sum_{j}^n I_j\cdot C_j + \sum_{j}^m W_j\cdot C_{n+j} \right)\\
= & (2(6x+10)+6(7x+4))\cdot(3(6x+10)+4(7x+4))-(11(7x+4)+6(6x+10)) \\
= & ((12x+7)+(3x+11))\cdot((5x+4)+(2x+3))-((12x+5)+(10x+8)) \\
= & (2x+5)\cdot(7x+7)-(9x) \\
= & (x^{2}+2\cdot7x+5\cdot7x+5\cdot7)-(9x) \\
= & (x^{2}+x+9x+9)-(9x) \\
= & x^{2}+x+9
\end{align*}
Учитывая экземпляр $I_1=11$, доказывающий предоставляет многочлен $x^2+x+9$ как доказательство. Чтобы верифицировать это доказательство, любой верификатор может найти целевой многочлен $T$ из QAP и разделить $P_{(I;W)}$ на $T$. В этом конкретном примере $P_{(I;W)}$ равен целевому многочлену $T$, и следовательно, он делится на $T$ с $P/T=1$. Верификатор таким образом верифицирует доказательство.
\begin{sagecommandline}
sage: F13 = GF(13)
sage: F13t.<t> = F13[]
sage: T = F13t(t^2 + t + 9)
sage: P = F13t((2*(6*t+10)+6*(7*t+4))*(3*(6*t+10)+4*(7*t +4))-(11*(7*t+4)+6*(6*t+10)))
sage: P == T
sage: P % T # remainder
\end{sagecommandline}

Чтобы привести пример ложного доказательства, рассмотрим строку $(<I_1>;<W_1,W_2,W_3,W_4>)=(<11>, <2, 3, 4, 8>)$. Выполняя схему, мы можем увидеть, что это не является корректным присваиванием и не решением R1CS, а следовательно, не конструктивным доказательством знания в $L_{3.fac\_zk}$. Однако доказывающий может использовать эти значения, чтобы сконструировать ложное доказательство $P_{(I;W)}$:
\begin{align*}
P'_{(I;W)}  = & \scriptstyle \left(A_0 + \sum_{j}^n I_j\cdot A_j + \sum_{j}^m W_j\cdot A_{n+j} \right) \cdot \left(B_0 + \sum_{j}^n I_j\cdot B_j + \sum_{j}^m W_j\cdot B_{n+j} \right) 
-\left(C_0 + \sum_{j}^n I_j\cdot C_j + \sum_{j}^m W_j\cdot C_{n+j} \right)\\
= & (2(6x+10)+8(7x+4))\cdot(3(6x+10)+4(7x+4))-(8(6x+10)+11(7x+4)) \\
= & 8x^{2}+6
\end{align*}
Учитывая экземпляр $I_1=11$, доказывающий поэтому предоставляет многочлен $8x^2+6$ как доказательство. Чтобы верифицировать это доказательство, любой верификатор может найти целевой многочлен $T$ из QAP и разделить $P_{(I;W)}$ на $T$. Однако деление многочленов имеет следующий остаток:
$$
(8x^{2}+6)/(x^{2}+x+9) =8+\frac{5x+12}{x^{2}+x+9} 
$$
Это подразумевает, что $P_{(I;W)}$ не делится на $T$, и следовательно, верификатор не верифицирует доказательство. Любой верификатор поэтому может показать, что доказательство ложно.
\begin{sagecommandline}
sage: F13 = GF(13)
sage: F13t.<t> = F13[]
sage: T = F13t(t^2 + t + 9)
sage: P = F13t((2*(6*t+10)+8*(7*t+4))*(3*(6*t+10)+4*(7*t+4))-(8*(6*t+10)+11*(7*t+4)))
sage: P == F13t(8*t^2 + 6)
sage: P % T # remainder
\end{sagecommandline}

\end{example}
